% ============================================================
%  第三卷 第十六章 核结构（§116–§122 + 习题）
%  底本：高教社中文版第6版（严肃 译，喀兴林 校）
%  OCR 底稿 full_merged.md 行 16327–17262
%  本章无插图（OCR 区间内无 images 行；任务单所提图44 经核 hash
%  a5ac31d165bb 与原书页面 p-402，实属 §105 氨分子位形图，已排在 ch13）
%
%  已修正的 OCR/原书错误（均对照 原书转换/pages/p-444…p-471.png 核实）：
%   [§116] OCR 章标题两行（"# 第十六章"+"# 核结构"）合并为 \chapter{核结构}；
%         "自旋都是 1/2"（OCR 分数错位）；"粒子自旋为 1/2 的多粒子系统"
%         "取 ±1/2 中的哪个值"（OCR 碎片重组，衍字 σ 删）；
%         行16347 "同位旋空间"后的乱码"$5,7,5$"按 p445 恢复为 ξ,η,ζ，
%         重复圈码去一（圈码3 置于"同位旋"之后）；
%         行16355/16357 补 T_ζ、ψ 对称性句按 p446 重组；
%         "对 ξηζ 空间内的转动"（OCR sinζ）；
%         τ̂ 下标 ξ,η,ζ 与 ŝ 下标 x,y,z 按 p446 恢复；τ̂_±、½+τ̂_ζ 矩阵
%         按 p446 重建；(116.2) 下标 对称/反对称 按 p447；
%         偶极矩例 f_p=er（OCR f_p 为 er 而 f_n=0 碎片重组）；
%   [§117] (117.1)(117.2) 下标 普通/交换 按 p448 恢复（OCR ∥∥/xx）；
%         习题1 Û_{ST}（OCR U_{sr}）、U_1/U_5（OCR U_t/U_s）；
%         习题2 下标 核（OCR *）：S_核,T_核；A^{5/3}（OCR A^{s/3}）；
%         (117.4) 双行合排；(117.2) 大括号按 p448；
%   [§118] (118.1) 箭头恢复；"核子速度 v_i"（OCR σ）；
%         ħl（OCR hl）；V̂_{sl}（OCR V_{,l}）；f=ħφ/rm_p；
%         (118.5) l·s 本征值两行右括号按 p454 重建；
%         (118.7) 下标原书印作俄文缩写 гр（基态），按文意改排"基态"；
%         (118.8) 壳层填充组表（OCR 成分式乱码）按 p456 整表重建；
%         幻数行"2,8,20,50,82,126."按 p456；
%         (118.9) OCR 双帽 \hat{\hat{μ}} 按 p456 改单帽；μ̄_z（OCR \vec）；
%         g_l=1, g_s=5.585; g_l=0, g_s=−3.826（OCR 乱码）恢复；
%         g_jĵ=g_sŝ+g_l l̂ 求平均式按 p457（(l̂−ŝ) 上带横线）；
%         行16613 句末衍字"1"删；(118.14) 2cħ（OCR ch）、下标 附加；
%         习题3 磁矩求和式：9/5→9/15、3/4→4/3（原书排印误，两项与末式
%         (1/15)(13μ_p+2μ_n) 不符，按逐项计算改正）；
%         习题5 (2) 式中间项 (ŝ·l̂)l̂−(l̂·ŝ)l̂ 原书印作减号，两项恒等相消、
%         公式不成立，按 §29 例题公式改正为加号（改正后恰可推出 (3) 式）；
%   [§119] 行16863 句末衍字"DY"删；双幻核 ^{208}_{82}Pb 按 p461 补上下标；
%         Ω=|ω_p±ω_n| 按 p462；0_g^±、0_u^∓ 态句按 p462 恢复（OCR 乱码）；
%         转动能级间距比例 1:3,3:7:12 中的逗号系俄文小数逗号遗留，
%         按 6:20:42:72=1:3.3:7:12 改为小数点；
%         §118 习题5 讨论段 OCR 漏句"在原子核的轴对称场中并不存在核子的
%         角动量矢量，把 σ̂ 看作类似于自旋为 1/2 算"按 p462 补回（实属
%         §119 正文，OCR 错位脱落）；
%         (119.3) (−1)^{J−1/2} 按 p463；(119.4)–(119.6) g_r（OCR g_t/g_s）；
%         (119.6) 前之 n̄（OCR n̂）按 p464 带横线照排；
%         习题1 n_i\overline{J}_i=Ω 按 p464；
%   [§120] (120.1)(120.2) 按 p465 恢复（p̂_i、平均横线）；ψ² 出积分号按 p466；
%         "关系为 √Z""与 R²Z² 成正比"按 p466 照排（§71 准经典估计 ψ(0)∝√Z）；
%   [§121] μ_核、μ_电子 下标恢复（OCR */♯）；(121.2) ħ/mc²、Rv_核/r³；
%         ĵ=−2μ_B c curl(ψ²ŝ)（OCR −zμ_B）；
%         习题1 "因而可写出"（OCR 乱码 κ̄）；(ŝ·n̄)(n̄·ĵ) 带横线 n 按 p468；
%         末句"正比于原子序"（OCR 句末衍 √Z 删）；
%   [§122] i_ζ²（OCR i_5）；(122.4) b̄_{ik}、(122.5)(122.6) 平均横线按 p470；
%         转动角速度 Ω（OCR 花体 𝒜）按 p471 恢复。
%   [全章] \mathrm{d} 收紧；函数名正体；μ_B、MeV、g·cm/s 单位正体；
%         粒子符号 p/n 及轨道字母 s,p,d,f,g,h,i 正体；半角标点改全角；
%         正文圈码共 27 处原样保留，脚注文本由脚注代理后续转录。
%
%  遗留问题：
%   - §118 习题5 (3) 式 j=l−1/2 情形：按 (2) 式（改正后）推导得
%     ±f̄(m_pR²/ħ²)(2j+1)/(4j)，原书/中译均印 (2j+1)/(4(j+1))，照排并存疑；
%   - (118.7) 下标"基态"系俄文 гр 之汉化，如需严格影印可回改 гр；
%   - §118 正文"见 §117 第三个附注"原书如此（p453，指速度关联项之脚注），
%     圈码编号系按页重计，照排；
%   - §119 (119.6) 前之 n̄（横线）与 §121 习题1 之 n̄ 按页面照排，
%     系"对体内运动平均后的 n"，与 (118.14) 中之 n 含义有别。
%
%  脚注转录（脚注代理完成，正文圈码 29 处全部替换为 \footnote，前后一致。
%  排版代理原记"27 处"系漏计：§116 行16341 跨页合并行含两处圈码、
%  §117 (s1·s2)(n·s1)(n·s2) 后的圈码2 在 OCR 中误作"(2)"、§118 "质心不动"
%  后的圈码1 在 OCR 中作 \textcircled 均未入账，实为 29 处）。
%  页面：原书转换/pages300/p-444…p-471；定位依据 脚注页码对照.txt 行 16337–17245。
%   p444 注1-2 —— §116 镜像核；isotopic/isobaric invariance 译名；
%   p445 注1-4 —— §116 pp/pn 散射实验（Breit-Condon-Present 1936）；质量差别
%        的电磁起源；同位旋创用者（Heisenberg 1932, Cassen-Condon 1936）；
%        文献中相反的指定；
%   p447 注1 —— §116 习题2（1−T̂² 算符形式见 §62 例题）；
%   p448 注1-2 —— §117 核子自旋-自旋作用非纯相对论性；核力不含赝标量假设；
%   p449 注1-2 —— §117 速度相干作用算符 [φ1+φ2P̂](L̂·Ŝ)（含公式）；氘核实验
%        对 T=0,S=1 与 T=1,S=0 作用的结论；
%   p450 注1 —— §117 饱和性（密度 n、总动能 ∼A·A^(2/3)，注内含公式）；
%   p452 注1 —— §118 原子中质心静止（本注亦为 p467 注1 之"见 452 页注"所指）；
%   p454 注1-2 —— §118 最轻核接近 LS 耦合；主量子数 n 从 l+1 起的对比；
%   p456 注1-2 —— §118 1f_(7/2) 与幻数 28；双幻核清单（注内含
%        ^4_2He_2、^16_8O_8、^40_20Ca_20、^208_82Pb_126、^4He）；
%   p457 注1 —— §118 "单粒子"磁矩计算只给上下限；
%   p458 注1 —— §118 习题3（Li^7 组态）；
%   p461 注1-2 —— §119 分类相似性指双原子分子非对称陀螺；Ω⩾0 定义
%        （Λ+Σ，Σ 可正可负）；
%   p462 注1 —— §119 球形核中 l 可定义（宇称与角动量同时守恒）；
%   p463 注1-2 —— §119 Ω=1/2 情形特点与 Λ 双线（注1 含两段）；μ'n+
%        g_r(J−Ωn) 写法仅适用于 Ω≠1/2；
%   p465 注1-2 —— §120 质心系总动能（注1 含公式 p_核²/2M+…）；计算忽略
%        相对论效应的条件 Ze²/ħc≪1；
%   p467 注1 —— §121 Fermi 计算条件 Ze²/ħc≪1（交叉引用"452 页注1"
%        用 {\symbol{"2460}} 排出圈码，避免源文件残留裸圈码字符）；
%   p468 注1 —— §121 见《场论》(43.7)，矢量 R 反方向；
%   p470 注1 —— §122 J_ξJ_ζ 等平均值为零之由（ψ_Jk 中 k 差为偶数，
%        见 §103）。
% ============================================================
\chapter{核结构}
\emergencystretch 3em

\section{同位旋不变性}

目前还没有完整的核力理论，核力是作用在核粒子（核子）之间并把它们束缚在原子核内的一种力.由于还没有一个完整的核力理论，所以对核力的描述在更大程度上需要依靠实验.

核子有两种，两者的主要区别在于它们的电性.质子（p）带有正电荷，而中子（n）是电中性的.它们的自旋都是$\frac{1}{2}$，质量几乎相等（分别为1836.1和1838.6电子质量）.这种相似性不是偶然的，撇开电性的差别，质子和中子是两个极为相似的粒子，这种相似性在本质上具有重要的意义.

业已发现，除了较弱的电力外，两个质子间的相互作用力非常相似于两个中子间的相互作用力，这称为核力的电荷对称性\footnote{它特别是表现在镜像核性质（结合能，能谱等等）的相似性方面.一对镜像核是指质子数和中子数相互对换的两个核.}.

在保持这种对称性的范围内，我们可以说，双质子（pp）系统和双中子（nn）系统具有性质上相同的态.当然，在这里重要的是，质子和中子服从同一种统计（即费米统计），因此pp和nn系统只允许有波函数$\psi(\boldsymbol{r}_1,\boldsymbol{\sigma}_1;\boldsymbol{r}_2,\boldsymbol{\sigma}_2)$对称性相同的态，即对粒子坐标和自旋的同时交换保持反对称性的那种态.

可是，电荷对称性不过是质子和中子间更深入的物理相似性即所谓同位旋不变性\footnote{同位旋不变性，英文为 isotopic invariance，也有称为 isobaric invariance.}的表现之一.这个性质不但导致pp和nn系统（可通过所有质子与所有中子的对换而相互得到）的相似性，而且还导致由不同粒子组成的pn系统与以上两个系统的相似性，当然，它们不可能是完全的相似，因为pn系统中的粒子并不全同，它的态肯定不限于波函数为反对称的态.但是我们发现，在pn系统的各种可能态中有些态在性质上和两个全同核子系统的态差不多精确相同\footnote{这可根据pp散射和pn散射的实验数据分析中得出（G.~Breit, E.~U.~Condon, R.~D.~Present, 1936）.}.这些态当然是由反对称波函数描写的（pn系统剩下的态由对称波函数所描写，在pp和nn系统中并不存在）.

同位旋不变性和电荷对称性一样，仅当略去电磁作用后才能成立.同位旋不变性为什么只是近似正确的另一个理由是，中子和质子间存在着很小的质量差别；如果中子和质子间真是精确对称的，它们的质量当然也应等同\footnote{实际上，中子和质子间的这个质量差别很可能还是由电磁原因造成的.}.

有一种方便的表述方式可以用来描述同位旋不变性.它是根据下列事实自然地得到的，即同位旋不变性相当于把核子系统的态按其坐标-自旋波函数$\psi$的对称性进行分类的可能性，而与所述及的核子类型无关，因此在所找的表述方式中，我们一定有可能定义一个描述系统状态的新量子数，用以唯一地确定函数$\psi$的对称性.与此类似的一种情况，我们早已在粒子自旋为$\frac{1}{2}$的多粒子系统中碰到过.我们在§63中看到，如果指定了这个系统的总自旋$S$，则其坐标波函数$\varphi$的对称性就被唯一地确定，而不管每个粒子的自旋分量$\sigma$究竟取$\pm\frac{1}{2}$中的哪个值.

因此在对同位旋不变性进行形式上的描述时，我们就有理由把中子和质子看作同一个粒子（核子）的两个不同“电荷态”，它们的区别在于一个新矢量$\boldsymbol{\tau}$的分量具有不同的值，这个新矢量$\boldsymbol{\tau}$在形式上与自旋为$\frac{1}{2}$的自旋矢量具有完全类似的性质.这个新量通常称为同位自旋或同位旋\footnote{先由 W.~Heisenberg（1932）所采用，并由 B.~Cassen 和 E.~U.~Condon（1936）拿来描写同位旋不变性.}，它是“同位旋空间”$\xi,\eta,\zeta$（当然它和实际空间毫无关系）中的一个矢量.

一个核子的同位旋的$\zeta$轴分量只能取$\tau_\zeta=\pm\frac{1}{2}$两个值.$+\frac{1}{2}$值被任意地指定给质子，$-\frac{1}{2}$值则指定给中子\footnote{文献中也能找到相反的指定.}.几个核子的同位旋可按通常自旋的相加法则相加成为系统的总同位旋.系统总同位旋的$\zeta$分量等于各个粒子的$\tau_\zeta$值之和.对于一个质子数为$Z$（即原子序数）中子数为$N$质量数$A=Z+N$的核，我们有
\begin{equation}
  T_\zeta=\sum\tau_\zeta=\frac{1}{2}(Z-N)=Z-\frac{1}{2}A.
\end{equation}
亦即当核子数固定后，$T_\zeta$能给出该系统的总电荷.因此$T_\zeta$显然是一个严格守恒量，它简单地表达了电荷的守恒.

正如总自旋$S$确定了自旋波函数的对称性那样，系统总同位旋的绝对值$T$确定了该系统“电荷部分”波函数$\omega$的对称性.从而也确定了坐标-自旋（即通常的）波函数$\psi$的对称性，这是因为核子系统的总波函数（即乘积$\psi\omega$）必须具备一定的对称性，它和所有的费米子系统一样，对两个粒子的坐标，自旋以及“电荷变量”$\tau_\zeta$的同时对换必须反对称.因此任一核子系统中波函数$\psi$存在着确定的对称性，在上述处理中被表达成为$T$的守恒性.

我们也可以换一种说法，所谓同位旋不变性，是指系统的性质对同位旋空间中的转动具有不变性.仅仅是$T_\zeta$值不同的各个态（$T$和其它量子数具有定值）具有相同的性质.电荷对称性（中子和质子互换后系统性质的不变性）不过是同位旋不变性的一个特例，可以看成是对所有$\tau_\zeta$的同时变号所具有的不变性，亦即绕同位旋空间$\xi\eta$平面内的一个轴旋转$180^\circ$所具有的不变性.

根据以上的处理显然可知，同位旋不变性必然会被库仑作用所破坏，库仑作用与电荷有关，也就是和同位旋的$\zeta$分量有关，故对$\xi\eta\zeta$空间内的转动并不具备不变性.

我们以双核子系统为例，它的总同位旋可取$T=1$和$T=0$两个值.$T=1$时$T_\zeta$分量的可能值为$1,0,-1$.根据(116.1)，相应的电荷值为$2,1,0$，亦即$T=1$的系统可以是pp，pn或nn.$T=1$的波函数电荷部分$\omega$是对称的（正如对称的自旋波函数对应于自旋$S=1$一样，参考§62）.所以$T=1$的态具有反对称的通常波函数$\psi$.$T=0$时只能有$T_\zeta=0$，相应的波函数$\omega$是反对称的；因此它只和pn系统中$\psi$波函数为对称的态有关.

同位旋对应于一个算符$\hat{\boldsymbol{\tau}}$，它作用在波函数的电荷变量$\tau_\zeta$上，正和自旋算符$\hat{\boldsymbol{s}}$作用在自旋变量$\boldsymbol{\sigma}$上一样.鉴于两者在形式上完全类似，$\hat{\tau}_\xi,\hat{\tau}_\eta,\hat{\tau}_\zeta$算符和$\hat{s}_x,\hat{s}_y,\hat{s}_z$算符一样，由(55.7)式的泡利矩阵给出.

下面来给出这些算符的某些组合形式，它们具有简单的直观含义.下列算符
\begin{equation*}
  \hat{\tau}_+=\hat{\tau}_\xi+\mathrm{i}\hat{\tau}_\eta
  =\begin{pmatrix}0&1\\0&0\end{pmatrix}
\end{equation*}
作用在中子波函数上把它变成质子波函数，作用在质子波函数上则等于零.同理，算符
\begin{equation*}
  \hat{\tau}_-=\hat{\tau}_\xi-\mathrm{i}\hat{\tau}_\eta
  =\begin{pmatrix}0&0\\1&0\end{pmatrix}
\end{equation*}
把质子变成中子，并把中子“湮没”掉.最后，下列算符
\begin{equation*}
  \frac{1}{2}+\hat{\tau}_\zeta=\begin{pmatrix}1&0\\0&0\end{pmatrix}
\end{equation*}
使质子波函数不变并把中子波函数湮没掉.上式乘$e$后可以称为核子的电荷算符.

现在来证明，两个粒子的对换算符$\hat{P}$可以用它们的同位旋算符$\hat{\boldsymbol{\tau}}_1,\hat{\boldsymbol{\tau}}_2$表达出来.根据定义，对换算符作用在双粒子系统波函数$\psi(\boldsymbol{r}_1,\boldsymbol{\sigma}_1;\boldsymbol{r}_2,\boldsymbol{\sigma}_2)$上的结果是使这两个粒子的坐标和自旋同时对换，亦即变量$\boldsymbol{r}_1,\boldsymbol{\sigma}_1$和$\boldsymbol{r}_2,\boldsymbol{\sigma}_2$对换.这个算符的本征值为$\pm1$，当它作用在对称的或反对称的$\psi$波函数上时就有
\begin{equation}
  \hat{P}\psi_{\text{对称}}=\psi_{\text{对称}};\qquad
  \hat{P}\psi_{\text{反对称}}=-\psi_{\text{反对称}}.
\end{equation}
前面已经讲过，$\psi_{\text{对称}}$和$\psi_{\text{反对称}}$分别对应于总同位旋$T=0$和$T=1$的电荷波函数$\omega_T$.由此可见，为了把$\hat{P}$表成作用于电荷变量上的算符，它就应该具有下列性质
\begin{equation}
  \hat{P}\omega_0=\omega_0,\qquad\hat{P}\omega_1=-\omega_1.
\end{equation}
这两个条件能被算符$1-\hat{\boldsymbol{T}}^2$所满足，这是很易看出的，只要注意到$\omega_T$是算符$\hat{\boldsymbol{T}}^2$的本征函数，具有本征值$T(T+1)$.最后，写出$\boldsymbol{T}=\boldsymbol{\tau}_1+\boldsymbol{\tau}_2$并且考虑到$\hat{\boldsymbol{\tau}}_1^2$和$\hat{\boldsymbol{\tau}}_2^2$都具有定值$\tau(\tau+1)=\frac{3}{4}$，即得所求的表式\footnote{这种形式的算符，已在§62例题中用粒子的通常自旋导出过.}
\begin{equation}
  \hat{P}=1-\hat{\boldsymbol{T}}^2
  =-\frac{1}{2}-2\hat{\boldsymbol{\tau}}_1\cdot\hat{\boldsymbol{\tau}}_2
\end{equation}
对于核子系统中各种物理量的矩阵元，存在着一定的同位旋选择定则（L.~A.~Radicati,1952）.设$F$为具有相加性的某个量（任意秩张量），亦即它对于整个系统而言的值等于对各个个别核子而言的值之和.我们把这个量的算符写成
\begin{equation*}
  \hat{F}=\sum_p\hat{f}_p+\sum_n\hat{f}_n,
\end{equation*}
式中分别对该系统内所有的质子和中子求和.这个表式可写成下列等同形式：
\begin{equation}
\begin{aligned}
  \hat{F}&=\sum\left(\frac{1}{2}+\hat{\tau}_\zeta\right)\hat{f}_p
  +\sum\left(\frac{1}{2}-\hat{\tau}_\zeta\right)\hat{f}_n=\\
  &=\frac{1}{2}\sum\left(\hat{f}_p+\hat{f}_n\right)
  +\sum\left(\hat{f}_p-\hat{f}_n\right)\hat{\tau}_\zeta,
\end{aligned}
\end{equation}
式中每项都是对所有核子（质子和中子）求和.(116.5)中的第一项是标量，第二项是同位旋空间中矢量的$\zeta$分量.因此它们对同位旋而言的选择定则，与普通空间中的标量和矢量对轨道角动量而言的选择定则（见§29）相同：同位旋标量只允许不改变$T$值的跃迁；同位旋矢量的$\zeta$分量只能有$\Delta T=0$或$\pm1$的跃迁矩阵元.对于$T_\zeta=0$的两个态，也就是对于中子数和质子数相等的系统，不能有$\Delta T=0$的跃迁，这是因为$\Delta T=0$的跃迁矩阵元是和$T_\zeta$成正比的（见(29.7)式）.

以原子核的偶极矩为例，$\boldsymbol{f}_p$为$e\boldsymbol{r}$而$\boldsymbol{f}_n=0$.(116.5)中的第一项为
\begin{equation*}
  \frac{1}{2}e\sum\boldsymbol{r}=\frac{e}{2m}\sum m\boldsymbol{r}.
\end{equation*}
因而和质心的径矢成正比，适当选取原点后可使它等于零.可见核偶极矩可化成同位旋矢量的$\zeta$分量.

\section{核力}

作用于核子之间的核力的主要特征是力程很短：在数量级为$10^{-13}\ \mathrm{cm}$的距离处指数式地衰减.

在非相对论极限下我们可以这样说，核力与核子的速度无关并且有势；核内的核子速度大约是光速的1/4（见后）.两核子的相互作用势能$U$不但和距离$r$有关而且相当强烈地和它们的自旋有关\footnote{从这一点讲来，粒子作用和电子作用有很大的不同，后者的自旋-自旋作用是纯相对论性的而且（在原子中）很小.}.$U$和$r$的确切关系当然只能由核力理论去确立，至于它和自旋的关系则可根据自旋算符的性质经过简单考虑后得出.

可供我们支配的与相互作用能量$U$有关的矢量一起只有三个，即两核子间的单位径矢量$\boldsymbol{n}$，以及两个核子的自旋$\boldsymbol{s}_1$和$\boldsymbol{s}_2$.根据自旋$\frac{1}{2}$算符的一般性质，它的任意函数可以化成线性函数（§55）.再考虑到乘积$\boldsymbol{n}\cdot\boldsymbol{s}$不是真标量而是赝标量（因为$\boldsymbol{n}$是极矢量而$\boldsymbol{s}$是轴矢量）.于是很明显，从$\boldsymbol{n},\boldsymbol{s}_1,\boldsymbol{s}_2$三个矢量出发只能构造出两个与自旋呈线性关系的独立标量，即$\left(\boldsymbol{s}_1\cdot\boldsymbol{s}_2\right)$\ ($\boldsymbol{n}\cdot\boldsymbol{s}_1)(\boldsymbol{n}\cdot\boldsymbol{s}_2)$\footnote{这里假定了核力是空间反演不变的，亦即不能含有赝标量.迄今为止没有实验否定这个假设.}.

从而，两核子的相互作用算符计及对自旋的关系后，可以写成三个独立项之和：
\begin{equation}
\begin{aligned}
  \hat{U}_{\text{普通}}&=U_1(r)+U_2(r)\left(\hat{\boldsymbol{s}}_1\cdot\hat{\boldsymbol{s}}_2\right)\,+\\
  &\quad+U_3(r)\left[3(\hat{\boldsymbol{s}}_1\cdot\boldsymbol{n})(\hat{\boldsymbol{s}}_2\cdot\boldsymbol{n})-\hat{\boldsymbol{s}}_1\cdot\hat{\boldsymbol{s}}_2\right],
\end{aligned}
\end{equation}
其中两项与自旋有关，一项无关.第三项被写成这样的形式，是使它对$\boldsymbol{n}$的各个方向平均后等于零.这一项所描述的力通常称为张量力.

(117.1)中的下标“普通”表明这个算符不影响核子的电荷态.实际上还存在着作用后使质子变成中子以及中子变成质子的相互作用.这种“交换”作用的算符与(117.1)不同之处在于还含有(116.4)的粒子对换算符：
\begin{equation}
  \hat{U}_{\text{交换}}=\left\{
  \begin{aligned}
    &U_4(r)+U_5(r)\left(\hat{\boldsymbol{s}}_1\cdot\hat{\boldsymbol{s}}_2\right)\,+\\
    &+U_6(r)\left[3(\hat{\boldsymbol{s}}_1\cdot\boldsymbol{n})(\hat{\boldsymbol{s}}_2\cdot\boldsymbol{n})-\hat{\boldsymbol{s}}_1\cdot\hat{\boldsymbol{s}}_2\right]
  \end{aligned}\right\}\hat{P}.
\end{equation}
总的相互作用算符为
\begin{equation}
  \hat{U}=\hat{U}_{\text{普通}}+\hat{U}_{\text{交换}}.
\end{equation}
可见两核子的相互作用要用六个不同的距离函数来描写.这些项一般讲来都属于同一数量级\footnote{还可提一下依赖于核子速度的相互作用.在速度的线性近似下，可以用一个具有$[\varphi_1(r)+\varphi_2(r)\hat{P}](\hat{\boldsymbol{L}}\cdot\hat{\boldsymbol{S}})$形式的算符来描写，其中$\boldsymbol{L}=\boldsymbol{r}\times\boldsymbol{p}$是两核子相对运动的轨道角动量，$\boldsymbol{p}$是动量，$\boldsymbol{S}=\boldsymbol{s}_1+\boldsymbol{s}_2$.这个算符含有两个$r$的函数.根据宇称及时间反演不变性，$\boldsymbol{p}\cdot\boldsymbol{n}$和$\boldsymbol{S}\cdot\boldsymbol{n}$等项被排除.}，

(117.1)和(117.2)中的自旋算符可以用总自旋算符$\hat{\boldsymbol{S}}$表出，把$\hat{\boldsymbol{S}}=\hat{\boldsymbol{s}}_1+\hat{\boldsymbol{s}}_2$和$\hat{\boldsymbol{S}}\cdot\boldsymbol{n}=\hat{\boldsymbol{s}}_1\cdot\boldsymbol{n}+\hat{\boldsymbol{s}}_2\cdot\boldsymbol{n}$平方起来并利用$\hat{\boldsymbol{s}}_1^2=\hat{\boldsymbol{s}}_2^2=\frac{3}{4}$，$(\hat{\boldsymbol{s}}_1\cdot\boldsymbol{n})^2=(\hat{\boldsymbol{s}}_2\cdot\boldsymbol{n})^2=\frac{1}{4}$（见(55.10)式），我们得
\begin{equation}
\begin{gathered}
  \hat{\boldsymbol{s}}_1\cdot\hat{\boldsymbol{s}}_2
  =\frac{1}{2}\left(\hat{\boldsymbol{S}}^2-\frac{3}{2}\right),\\
  \left(\hat{\boldsymbol{s}}_1\cdot\boldsymbol{n}\right)
  \left(\hat{\boldsymbol{s}}_2\cdot\boldsymbol{n}\right)
  =\frac{1}{2}\left[\left(\hat{\boldsymbol{S}}\cdot\boldsymbol{n}\right)^2-\frac{1}{2}\right].
\end{gathered}
\end{equation}
算符$\hat{\boldsymbol{S}}^2$和$\hat{\boldsymbol{S}}$对易，因此对于(117.1)和(117.2)中前两项所代表的作用讲来，系统的总自旋矢量是守恒的.张量作用含有算符$(\hat{\boldsymbol{S}}\cdot\boldsymbol{n})^2$，它和$\hat{\boldsymbol{S}}^2$对易但和矢量$\hat{\boldsymbol{S}}$本身不对易.结果只有总自旋的绝对值守恒，它的方向不守恒.

双核子系统的总自旋$S$可取0和1，总同位旋$T$也是这样，因此这个系统所有的态可按$S$和$T$值的不同分成四类.每类的态具有相互作用算符$A(r)$（对$S=0$）或$A(r)+B(r)\left[\left(\hat{\boldsymbol{S}}\cdot\boldsymbol{n}\right)^2-\dfrac{2}{3}\right]$（对$S=1$），可按每类情况由一般表式(117.3)化出（见题1）\footnote{有关氘核性质的实验表明，$T=0,S=1$的核子作用具有强的引力和一个深的“势阱”（张量力的存在使它难于用函数$A(r)$和$B(r)$的性质表达出来）.此外，根据观测到的氘核四极矩的符号，可知这个态的张量力中$B(r)$的系数是负的.根据核子散射实验结果，可知$T=1,S=0$的作用也是引力作用，但比较弱，特别是不能形成双粒子的稳定系统.}.

$S$和$T$给定后，系统的态是按总角动量值$J$和宇称分类的.我们知道，$T=0$和$T=1$分别对应于波函数$\psi$为对称的和反对称的态.另一方面，$S$值确定了波函数对自旋变量的对称性（$S=1$对称，$S=0$反对称）.显然，当$S$和$T$给定以后，波函数对空间变量的对称性（态的宇称）也被确定.同位旋$T=0$的态只能是偶的三重态（$S=1$）或奇的单态（$S=0$）；而同位旋$T=1$的态均为奇的三重态或偶的单态.

由于自旋作为矢量并不守恒，轨道角动量一般也不守恒，只有两者之和$\boldsymbol{J}=\boldsymbol{L}+\boldsymbol{S}$是守恒的.可是$\boldsymbol{L}$的绝对值也有可能守恒，这是因为给定了$J,S$和宇称（或$J,S$和$T$）后，有可能只有一个特定的$L$值可以与之相容（记得双粒子系统的宇称为$(-1)^L$）.例如$S=1,J=1$的奇态只能有$L=1$，即${}^{3}\mathrm{P}_1$.在其它情形下，给定了$J,S$和宇称后可以有两个不同的$L$值，因而$L$是不守恒的.例如$S=1,J=2$的奇态可以有$L=1$或$L=3$，这是一个组合态${}^{3}\mathrm{P}_2+{}^{3}\mathrm{F}_2$.

因此我们得到双核子系统的下列各种可能态（符号$\pm$代表宇称）：
\begin{equation*}
\begin{array}{rl}
  T=1: & {}^{3}\mathrm{P}_0^-,\ {}^{3}\mathrm{P}_1^-,\ ({}^{3}\mathrm{P}_2+{}^{3}\mathrm{F}_2)^-,\ {}^{3}\mathrm{F}_3^-,\ \cdots\\
  & {}^{1}\mathrm{S}_0^+,\ {}^{1}\mathrm{D}_2^+,\ {}^{1}\mathrm{G}_4^+,\ \cdots\\
  T=0: & ({}^{3}\mathrm{S}_1+{}^{3}\mathrm{D}_1)^+,\ {}^{3}\mathrm{D}_2^+,\ ({}^{3}\mathrm{D}_3+{}^{3}\mathrm{G}_3)^+,\ \cdots\\
  & {}^{1}\mathrm{P}_1^-,\ {}^{1}\mathrm{F}_3^-,\ \cdots
\end{array}
\end{equation*}
核力一般讲来并不是相加性的.这就是说，两个以上核子所组成的系统中，它的核力作用并不等于其中所有各对粒子的核力作用之和.然而，与二体作用相比较，三体作用和多体作用看来并不重要，因此在讨论复杂核的性质时，在很大程度上我们仍可以二体作用的性质为基础.

原子核实验结果表明，当粒子数$A$增大时，核子系统开始类似于宏观的“核物质”，它的体积和能量都与$A$成正比地增长（质子间库仑作用以及核的自由表面所产生的那些效应都是很小的）.产生这种现象的核力性质称为饱和性.

饱和性的存在使核子的二体作用函数$U_1,\cdots,U_6$受到一定的限制.假定所有粒子都集中在核力作用半径那样大小的一个体积内，使得每对粒子间都有相互作用存在.如果有这样一种核子组态（以及这样一些自旋取向）其中每对间的作用力都是吸引力，那么这个系统的势能是负的并且与$A^2$成正比，其动能是正的并且和$A^{5/3}$成正比（$A$的较小幂次）\footnote{集中于某一给定体积内的粒子其密度$n$与粒子数$A$成正比，每个粒子的动能与$n^{2/3}$成正比（参考(70.1)），故总动能$\sim A\cdot A^{2/3}$.}.显然，在这种条件下会有足够多的核子集中到一个与$A$无关的小体积中去，即不能形成核物质.由此可见，核力饱和性可以表成这样的条件：与$A^2$成正比的负作用能量的那些组态都不存在（见题2）.

核物质体积与其粒子数的正比性可用下式表出：
\begin{equation}
  R=r_0A^{1/3}.
\end{equation}
上式给出了核半径$R$和核中粒子数$A$的关系.实验结果（电子和核的散射）给出$r_0=1.1\times10^{-13}\ \mathrm{cm}$.

我们可求出核物质中核子动量的极限值（参考§70）.物理空间单位体积内动量为$p\leqslant p_0$的粒子所占的相空间体积为$4\pi p_0^3/3$，除以$(2\pi\hbar)^3$后即得“相格”数，每格中可以同时占有两个质子和两个中子.令中子数等于质子数，我们得
\begin{equation*}
  4\,\frac{4\pi}{3}\left(\frac{p_0}{2\pi\hbar}\right)^3=\frac{A}{V},
\end{equation*}
$V$是原子核体积.把(117.5)代入得
\begin{equation*}
  p_0=\left(\frac{3\pi^2A}{2V}\right)^{\frac{1}{3}}\hbar
  =(9\pi)^{1/3}\frac{\hbar}{2r_0}
  =1.4\times10^{-14}\ \mathrm{g}\cdot\mathrm{cm/s}.
\end{equation*}
相应的动能$p_0^2/2m_\mathrm{p}\sim40\ \mathrm{MeV}$（$m_\mathrm{p}$为核子质量），而速度为
\begin{equation*}
  \frac{p_0}{m_\mathrm{p}}\approx\frac{c}{4}.
\end{equation*}

\begin{xiti}
\setcounter{equation}{0}\renewcommand{\theequation}{\arabic{equation}}%

\noindent{\heiti 习题1}\quad 对$S,T$具有定值的各种双核子态，求其相互作用算符.

解：根据一般表式(117.1)---(117.3)，应用(116.3)和(117.4)，得所求的$\hat{U}_{ST}$算符：
\begin{equation*}
\begin{aligned}
  &\hat{U}_{00}=U_1-\frac{3}{4}U_2+U_4-\frac{3}{4}U_5,\\
  &\hat{U}_{01}=U_1-\frac{3}{4}U_2-U_4+\frac{3}{4}U_5,\\
  &\hat{U}_{10}=U_1+\frac{1}{4}U_2+U_4+\frac{1}{4}U_5
  +\frac{1}{2}\left(U_3+U_6\right)\left[3\left(\hat{\boldsymbol{S}}\cdot\boldsymbol{n}\right)^2-2\right],\\
  &\hat{U}_{11}=U_1+\frac{1}{4}U_2-U_4-\frac{1}{4}U_5
  +\frac{1}{2}\left(U_3-U_6\right)\left[3\left(\hat{\boldsymbol{S}}\cdot\boldsymbol{n}\right)^2-2\right].
\end{aligned}
\end{equation*}

\noindent{\heiti 习题2}\quad 求核力饱和性的条件，假定张量力不存在，其它各型的力假定具有相等的作用半径.

解：对核子数为$A$的系统，考虑几种极端情形的态（其它各种情形都介于其间）.写出这个系统中一个“平均”核子对的相互作用能为正的各种条件.

假定原子核的总自旋和总同位旋都呈极大值：$S_{\text{核}}=T_{\text{核}}=\dfrac{1}{2}A$（当系统中的粒子都是质子而且它们的自旋都平行时）.则每对核子具有$S=T=1$，要写出的条件为
\begin{equation}
  U_{11}>0.
\end{equation}
其次，设$T_{\text{核}}=\dfrac{1}{2}A$，$S_{\text{核}}=0$.则每对核子的$T=1$，每个粒子的$s_z$平均值为零.后者表明核子的$s_z=\dfrac{1}{2}$和$s_z=-\dfrac{1}{2}$是等概率的.在这些条件下，一对核子处于$S=0$态或$S=1$态的概率分别为1/4和3/4（与$S_z$值的个数$2S+1$成正比）.因此核子对平均能量为正的条件为
\begin{equation}
  \frac{1}{4}U_{01}+\frac{3}{4}U_{11}>0.
\end{equation}
同样地，讨论$T_{\text{核}}=0$，$S_{\text{核}}=\dfrac{1}{2}A$的态，得条件
\begin{equation}
  \frac{1}{4}U_{10}+\frac{3}{4}U_{11}>0.
\end{equation}
在$T_{\text{核}}=S_{\text{核}}=0$的那些态中，核子对具有$S=T=1$的概率为$3/4\times3/4$，具有$T=1$，$S=0$的概率为$3/4\times1/4$，以此类推.从而得条件
\begin{equation}
  \frac{9}{16}U_{11}+\frac{3}{16}\left(U_{10}+U_{01}\right)+\frac{1}{16}U_{00}>0.
\end{equation}
最后，设系统由$\dfrac{1}{2}A$个质子和$\dfrac{1}{2}A$个中子所组成，所有质子的自旋平行并和所有中子的自旋反平行.单个核子为$\mathrm{p}\left(\tau_\zeta=\dfrac{1}{2}\right)$或$\mathrm{n}\left(\tau_\zeta=-\dfrac{1}{2}\right)$的概率相等，核子对具有$T=0$的概率为1/4，由于这一对核子中一个是p另一个是n，故$S_z=0$.这个$S_z$值以相等的概率来自$S=0$或$S=1$的态.因此核子对处于$T=0$，$S=0$态和$T=0$，$S=1$态的概率都等于$1/4\times1/2=1/8$.$T=1$，$S=0$的态也有这样的概率，剩下5/8的概率就属于$T=S=1$的态.因而得条件
\begin{equation}
  \frac{1}{8}\left(U_{00}+U_{01}+U_{10}\right)+\frac{5}{8}U_{11}>0.
\end{equation}
不等式(1)---(5)构成核力饱和性所需的条件.

\end{xiti}
\renewcommand{\theequation}{\arabic{section}.\arabic{equation}}%

\section{壳层模型}

原子核的许多性质可用壳层模型很好地描述，它和原子的电子壳层结构基本上类似.这个模型中，原子核内的每个核子可以看作是在所有其余核子所组成的自洽场中运动：由于核力的作用范围很小，这个场一超出核“表面”所围的体积就很快地衰减.与此相应，整个原子核的态可用指定各个单核子态的办法来描写.

自洽场是球对称的，对称中心当然就是原子核的质心，这样一来，就产生了下列困难.在自洽场法中，系统的波函数是由单核子波函数的乘积（或适当对称化后的乘积之和）构成的，可是这样的函数无法保持质心不动\footnote{对于原子中的电子不会产生这种困难，因为它的质心和不动的原子核位置重合，必然是静止的.}：由这种函数算出的质心平均速度虽然等于零，但速度值本身的概率并不等于零.

当我们用自洽场法的波函数$\psi(\boldsymbol{r}_1,\cdots,\boldsymbol{r}_A)$计算任一物理量时，可以采用先消除质心运动的办法来避免这一困难.设$f(\boldsymbol{r}_i,\boldsymbol{p}_i)$为某一物理量，它是核子坐标和动量的函数.当用函数组$\psi$计算它的矩阵元时，我们必须在不改变$\psi(\boldsymbol{r}_i)$的情况下把函数$f$中的宗量改为
\begin{equation}
  \boldsymbol{r}_i\to\boldsymbol{r}_i-\boldsymbol{R},\qquad
  \boldsymbol{p}_i\to\boldsymbol{p}_i-\frac{1}{A}\boldsymbol{P}.
\end{equation}
式中$\boldsymbol{R}$是原子核质心的径矢，$A$是核内粒子数，$\boldsymbol{P}$是整体运动的动量.(118.1)中的第二项相当于从核子速度$\boldsymbol{v}_i$中减去质心速度$\boldsymbol{V}$，$\boldsymbol{P}$和$\boldsymbol{V}$的关系为$\boldsymbol{P}=Am_\mathrm{p}\boldsymbol{V}$（S.~Gartenhaus, C.~Schwartz,1957）.

例如原子核偶极矩算符为$\boldsymbol{d}=e\sum\boldsymbol{r}_p$，式中对核内所有质子求和.用自洽场法算其矩阵元时，这个算符必须改成$e\sum(\boldsymbol{r}_p-\boldsymbol{R})$.原子核的质心坐标为
\begin{equation*}
  \boldsymbol{R}=\frac{1}{A}\left(\sum_p\boldsymbol{r}_p+\sum_n\boldsymbol{r}_n\right),
\end{equation*}
式中对所有质子和中子求和.由于核内质子数为$Z$，偶极矩算符最后改成
\begin{equation}
  e\sum_p\boldsymbol{r}_p\to e\left(1-\frac{Z}{A}\right)\sum_p\boldsymbol{r}_p
  -e\frac{Z}{A}\sum_n\boldsymbol{r}_n,
\end{equation}
上式中的质子具有“有效电荷”$e(1-Z/A)$，中子具有“电荷”$-eZ/A$.由(118.2)知，偶极矩改正项的相对数量级为1.不难求出，磁矩和更高级电多极矩改正项的相对数量级为$1/A$.

非相对论近似中，核子和自洽场的相互作用与该核子的自旋无关：这种关系只能和$\hat{\boldsymbol{s}}\cdot\boldsymbol{n}$成正比，$\boldsymbol{n}$是沿核子径矢$\boldsymbol{r}$的单位矢量，这个乘积是一个赝标量而不是真标量.

但当计及依赖于粒子速度的相对论项以后，核子能量就会依赖于自旋.其中的最大项与速度成正比.从$\boldsymbol{s}$,$\boldsymbol{n}$和$\boldsymbol{v}$三个矢量出发可组成一个真标量$\boldsymbol{n}\times\boldsymbol{v}\cdot\boldsymbol{s}$，因此原子核内核子的自旋-轨道耦合算符为
\begin{equation}
  \hat{V}_{sl}=-\varphi(r)\,\boldsymbol{n}\times\hat{\boldsymbol{v}}\cdot\hat{\boldsymbol{s}},
\end{equation}
$\varphi(r)$是$r$的某个函数，见§117第三个附注.由于$m_\mathrm{p}\boldsymbol{r}\times\boldsymbol{v}$是粒子的轨道角动量$\hbar\boldsymbol{l}$，(118.3)也可写成
\begin{equation}
  \hat{V}_{sl}=-f(r)\,\hat{\boldsymbol{l}}\cdot\hat{\boldsymbol{s}},
\end{equation}
其中$f=\hbar\varphi/rm_\mathrm{p}$.应该强调这个作用是$v/c$的一级效应，而原子中电子的自旋-轨道耦合是二级效应（§72），这个差别是由于核力即使在非相对论近似中也依赖于自旋，而电子的非相对论相互作用（库仑力）是与自旋无关的.

自旋-轨道作用能量主要集中在原子核表面附近，亦即函数$f(r)$在核内衰减.这是因为这种作用在无限的核物质中根本不会存在，只要考虑到此时的系统是均匀的，显然不再存在某种优先的$\boldsymbol{n}$方向.

作用项(118.4)把轨道角动量为$l$的核子能级分裂成为角动量$j=l\pm\dfrac{1}{2}$的两个能级.

由于（按(31.3)式）
\begin{equation}
  \left.
  \begin{aligned}
    \boldsymbol{l}\cdot\boldsymbol{s}&=\frac{1}{2}l,
    &\qquad& j=l+\frac{1}{2}\\
    &=-\frac{1}{2}(l+1),&& j=l-\frac{1}{2}
  \end{aligned}
  \right\},
\end{equation}
分裂量为
\begin{equation}
  \Delta E=E_{l-\frac{1}{2}}-E_{l+\frac{1}{2}}
  =\overline{f(r)}\left(l+\frac{1}{2}\right).
\end{equation}
实验表明，$j=l+\dfrac{1}{2}$的能级（$\boldsymbol{l}$和$\boldsymbol{s}$平行）低于$j=l-\dfrac{1}{2}$的能级，这意味着$f(r)>0$.

原子核中核子的自旋-轨道耦合要比该核子与自洽场的作用弱，但是一般讲来它要比原子核内两个核子的直接作用能量来得大，因后一种作用随着原子量的增加而更快地衰减.

各种相互作用能的上述大小关系，使核能级必须按$jj$耦合分类：各个核子的自旋和轨道角动量相加成为总角动量$\boldsymbol{j}=\boldsymbol{l}+\boldsymbol{s}$，由于$\boldsymbol{l}$和$\boldsymbol{s}$间的关系不受粒子间直接作用的影响（M.~G\"oppert-Mayer, 1949; O.~Haxel, J.~H.~D.~Jensen, H.~E.~Suess, 1949）$j$具有定值\footnote{只有对最轻的原子核，耦合才接近于$LS$.}.随后单个核子的$j$相加成为原子核的总角动量$\boldsymbol{J}$（$\boldsymbol{J}$通常简称为核自旋，犹如把原子核当作一个基本粒子）.从这方面看来，核能级的分类与原子能级根本不同：在原子的电子壳层中，相对论性的自旋-轨道耦合一般地小于直接的电作用和交换作用，故其能级分类通常以$LS$耦合为基础.

原子核中每个核子的态由它的角动量$j$和它的宇称来描写.尽管矢量$\boldsymbol{l}$和$\boldsymbol{s}$并不分别守恒，可是核子轨道角动量的绝对值可以具有定值.因为角动量$j$可以来自$l=j-\dfrac{1}{2}$的态或者来自$l=j+\dfrac{1}{2}$的态，对于给定的$j$（半整数），这两个态具有不同的宇称$(-1)^l$，所以当$j$和宇称都确定后，量子数$l$也就被确定了.

具有给定$l$和$j$值的各个核子态习惯上按“主量子数”$n$编号（按能量的递增次序），$n$从1开始取整数值\footnote{与原子中电子能级的通常做法不同，在那里$n$的取值是从$l+1$开始的.}.各种态记作$1\mathrm{s}_{1/2},1\mathrm{p}_{1/2},1\mathrm{p}_{3/2}$，等等.其中字母前的数字为主量子数，字母s，p，d，$\cdots$按惯例代表$l$值，下标为$j$值.具有给定$n,l,j$值的一个态中可以同时具有不超过$2j+1$个中子和不超过$2j+1$个质子.

整个核的态（组态为给定）按惯例用$J$值及该态的宇称符号$+$或$-$来标志（后者在壳层模型中由所有核子$l$值的代数和的奇偶来确定）.

根据有关核性质实验结果的分析，有可能导出有关核能级位置的一系列规则.首先，我们发现核子能量随轨道角动量$l$增大.产生这个规则的原因是，当$l$增大时粒子的离心能随之增大，从而使结合能减小.

其次，对于一个给定的$l$，$j=l+\dfrac{1}{2}$的能级（相当于矢量$\boldsymbol{l}$和$\boldsymbol{s}$平行）位于$j=l-\dfrac{1}{2}$能级之下.这个规则已经在前面提及，它和原子核内核子的自旋-轨道耦合性质有关.

下面的规则与原子核的同位旋有关，已知同位旋的分量$T_\zeta$是由该核的质量数和质子数确定的（见(116.1)式）.对于一个给定的$T_\zeta$值，同位旋的绝对值可取$T\geqslant|T_\zeta|$的任意值.一般讲来，原子核的基态具有这些同位旋允许值中的最小值，即
\begin{equation}
  T_{\text{基态}}=|T_\zeta|=\frac{1}{2}(N-Z).
\end{equation}
这个规则来自中子-质子相互作用的一个性质，即在np系统中同位旋$T=0$的态（即氘核态）要比$T=1$态的结合能大，见§117的倒数第二个脚注.

我们也能对基态核的自旋建立起一些规则.这些规则确定了单核子角动量$j$如何相加成为核的总自旋.它表现在原子核中处于相同态的质子或中子有以相反的角动量配合“成对”的趋势.这种pp和nn对的结合能约为1或2~MeV的数量级.

这个现象特别表现在偶-偶核中（原子核含有偶数个质子和偶数个中子），上述规则使核子角动量成对地抵消掉，结果使这种核的总角动量等于零.

但如原子核具有奇数个质子或中子，并且满壳层外的所有核子处于相同态中，该核的总角动量通常就等于一个核子的角动量，因为所有的质子和中子配对以后只剩下了一个核子（满壳层的总角动量必然为零）.

对于奇-奇核（$Z$和$N$都是奇数），没有一般规则足以确定基态的自旋.

对原子核内各个壳层具体填充方式的讨论，需要对现有各种实验数据进行详细的分析，这就超出了本书的范围.下面只想大致地提几点.

研究原子性质时我们曾经看到，电子态可以分成若干个组，每当填满一组转入下一组时，电子的结合能就下降.对原子核也有类似的情况，核子态可分成下列各组：
\begin{equation}
  \left.
  \begin{array}{ll}
    1\mathrm{s}_{1/2}
    & \begin{array}{c}\text{核子数}\\ 2\end{array}\\
    1\mathrm{p}_{3/2},\,1\mathrm{p}_{1/2}, & 6\\
    1\mathrm{d}_{5/2},\,1\mathrm{d}_{3/2},\,2\mathrm{s}_{1/2} & 12\\
    1\mathrm{f}_{7/2},\,2\mathrm{p}_{3/2},\,1\mathrm{f}_{5/2},\,2\mathrm{p}_{1/2},\,1\mathrm{g}_{9/2} & 30\\
    2\mathrm{d}_{5/2},\,1\mathrm{g}_{7/2},\,1\mathrm{h}_{11/2},\,2\mathrm{d}_{3/2},\,3\mathrm{s}_{1/2}, & 32\\
    2\mathrm{f}_{7/2},\,1\mathrm{h}_{9/2},\,1\mathrm{i}_{13/2},\,2\mathrm{f}_{5/2},\,3\mathrm{p}_{3/2},\,3\mathrm{p}_{1/2} & 44
  \end{array}
  \right\}
\end{equation}
每组给出了质子或中子的总空位数.根据这些数字，每当原子核中的质子总数$Z$或中子总数$N$等于下列各数之一时，就有一个组被填满.
\begin{equation*}
  2,8,20,50,82,126.
\end{equation*}
这些数通常称为幻数\footnote{$1\mathrm{f}_{7/2}$各态（8 个空位）有时单独编成一组，所以 28 也具有某种幻数性质.}.

$Z$和$N$都是幻数的“双幻”核特别稳定.与邻近的核相比，它们再结合一个核子的能力特别弱，从而它们的第一激发态特别高\footnote{这些核有${}_2^4\mathrm{He}_2$，${}_8^{16}\mathrm{O}_8$，${}_{20}^{40}\mathrm{Ca}_{20}$，${}_{82}^{208}\mathrm{Pb}_{126}$，$^4\mathrm{He}$核不能再加进一个核子.}.

(118.8)所列各组的态，大致上反映了一些核的填充次序.但在实际上，却发现填充过程是相当不规则的，此外还应注意到，在不接近于幻数的重核中，能级间距可能和“成对能量”差不多大小，单个核子对作为一个状态成分的态概念本身也在很大程度上失去意义.

我们对壳模型中核磁矩的计算作几点说明.我们所指的当然是对核内的粒子运动平均以后的磁矩.这个平均磁矩$\bar{\boldsymbol{\mu}}$显然是沿核自旋$\boldsymbol{J}$的方向，它是核内唯一的特殊方向，因此它的算符为
\begin{equation}
  \hat{\boldsymbol{\mu}}=\mu_0 g\hat{\boldsymbol{J}},
\end{equation}
$\mu_0$是核磁子，$g$是回转磁因子.磁矩分量的本征值为$\bar{\mu}_z=\mu_0 gM_J$.它的最大值$\mu=\mu_0 gJ$通常简称为核磁矩$\mu$（参考(111.1)式）.采用此记号后有
\begin{equation}
  \hat{\boldsymbol{\mu}}=\mu\hat{\boldsymbol{J}}/J.
\end{equation}
核磁矩是由满壳层以外的核子磁矩组合而成的，因为满壳层内核子角动量已经抵消掉.每个核子在核内产生的磁矩由两部分组成：自旋部分和（质子情形下）轨道部分，亦即可表成$g_s\hat{\boldsymbol{s}}+g_l\hat{\boldsymbol{l}}$（此后我们略去因子$\mu_0$，即核磁矩通常以核磁子为单位）.自旋的和轨道的旋磁因子，对质子为$g_l=1$，$g_s=5.585$；对中子为$g_l=0$，$g_s=-3.826$.

对核内的核子运动平均以后，磁矩就正比于$j$，把它写成$g_jj$形式后，我们有
\begin{equation*}
  g_j\hat{\boldsymbol{j}}=g_s\hat{\boldsymbol{s}}+g_l\hat{\boldsymbol{l}}
  =\frac{1}{2}\left(g_l+g_s\right)\hat{\boldsymbol{j}}
  +\frac{1}{2}\left(g_l-g_s\right)\overline{\left(\hat{\boldsymbol{l}}-\hat{\boldsymbol{s}}\right)}
\end{equation*}
上式两边乘$\hat{\boldsymbol{j}}=\hat{\boldsymbol{l}}+\hat{\boldsymbol{s}}$，取其本征值得
\begin{equation*}
\begin{aligned}
  g_jj(j+1)&=\frac{1}{2}\left(g_l+g_s\right)j(j+1)\,+\\
  &\quad+\frac{1}{2}\left(g_l-g_s\right)\left[l(l+1)-s(s+1)\right],
\end{aligned}
\end{equation*}
令$s=\dfrac{1}{2}$，$j=l\pm\dfrac{1}{2}$，得
\begin{equation}
  g_j=g_l\pm\frac{g_s-g_l}{2l+1},\qquad j=l\pm\frac{1}{2}.
\end{equation}
取上述旋磁因子值，对质子磁矩$\mu_\mathrm{p}=g_jj$有
\begin{equation}
  \left.
  \begin{aligned}
    j=l-\frac{1}{2},\quad&\mu_\mathrm{p}=\left(1-\frac{2.29}{j+1}\right)j,\\
    j=l+\frac{1}{2},\quad&\mu_\mathrm{p}=j+2.29.
  \end{aligned}
  \right\}
\end{equation}
对中子有
\begin{equation}
  \left.
  \begin{aligned}
    j=l-\frac{1}{2},\quad&\mu_\mathrm{n}=\frac{1.91}{j+1}j,\\
    j=l+\frac{1}{2},\quad&\mu_\mathrm{n}=-1.91.
  \end{aligned}
  \right\}
\end{equation}
（T.~Schmidt,1937）.

如果封闭壳层外只有一个核子，(118.12)和(118.13)直接给出了核磁矩.对于两个核子，磁矩的相加也是很简单的（见题1）.当核子数超过2时，磁矩的平均必须用该系统的波函数，后者由单个核子波函数以适当的方式构成.如果核子组态以及整个原子核的态已经给定，当所给组态中具有给定$J$和$T$值的态只有一个的时候，系统的波函数能够唯一地构成（例如见题3）.否则的话，原子核的态是几个（$J$和$T$相同的）独立态的混合，而核波函数中的线性组合系数一般讲来还是未知的\footnote{但是，对核磁矩的“单粒子”计算，实际上相当不精确.这时(118.12)和(118.13)中的两对值不是磁矩的精确值而只是其上下限.}.

最后可指出，原子核内核子自旋-轨道耦合的存在，使得核内质子产生一个(118.9)式以外的附加磁矩（M.~G\"opper-Mayer, J.~H.~D.~Jensen,1952）.其理由是，当有外场存在时，显含粒子速度的相互作用算符中应把动量$\hat{\boldsymbol{p}}$改成$\hat{\boldsymbol{p}}-e\boldsymbol{A}/c$.(118.3)式作此替代后，采用(111.7)式的矢势，我们发现质子哈密顿算符中含有下列附加项：
\begin{equation*}
\begin{aligned}
  \varphi(r)\frac{e}{cm_\mathrm{p}}\boldsymbol{n}\times\boldsymbol{A}\cdot\hat{\boldsymbol{s}}
  &=f(r)\frac{e}{2c\hbar}\boldsymbol{r}\times(\boldsymbol{H}\times\boldsymbol{r})\cdot\hat{\boldsymbol{s}}=\\
  &=f(r)\frac{e}{2c\hbar}\boldsymbol{r}\times(\hat{\boldsymbol{s}}\times\boldsymbol{r})\cdot\boldsymbol{H}.
\end{aligned}
\end{equation*}
这一项等价于附加磁矩的出现，其算符为
\begin{equation}
\begin{aligned}
  \hat{\boldsymbol{\mu}}_{\text{附加}}
  &=-\frac{e}{2c\hbar}f(r)\,\boldsymbol{r}\times(\hat{\boldsymbol{s}}\times\boldsymbol{r})=\\
  &=-\frac{e}{2c\hbar}r^2f(r)\left\{\hat{\boldsymbol{s}}-(\hat{\boldsymbol{s}}\cdot\boldsymbol{n})\boldsymbol{n}\right\}
\end{aligned}
\end{equation}

\begin{xiti}
\setcounter{equation}{0}\renewcommand{\theequation}{\arabic{equation}}%

\noindent{\heiti 习题1}\quad 试求双核子系统（具有总角动量$\boldsymbol{J}=\boldsymbol{j}_1+\boldsymbol{j}_2$）的磁矩，用两个核子磁矩$\boldsymbol{\mu}_1$和$\boldsymbol{\mu}_2$表出.

解：与推导(118.11)式类似，我们得
\begin{equation*}
  \frac{\mu}{j}=\frac{1}{2}\left(\frac{\mu_1}{j_1}+\frac{\mu_2}{j_2}\right)
  +\frac{1}{2}\left(\frac{\mu_1}{j_1}-\frac{\mu_2}{j_2}\right)
  \frac{\left(j_1-j_2\right)\left(j_1+j_2+1\right)}{J(J+1)}.
\end{equation*}

\noindent{\heiti 习题2}\quad 试求三核子系统的各种可能态，每个核子的角动量$j=3/2$（主量子数相同）.

解：与§67中求等效电子系统的各种可能态类似，每个核子可以处于$(m_j,\tau_\zeta)$值不同的下列八个态中的一个：
\begin{equation*}
\begin{aligned}
  &(3/2,1/2),\,(1/2,1/2),\,(-1/2,1/2),\,(-3/2,1/2),\\
  &(3/2,-1/2),\,(1/2,-1/2),\,(-1/2,-1/2),\,(-3/2,-1/2).
\end{aligned}
\end{equation*}
把其中的三个不同态组合起来，可得$(M_J,T_\zeta)$值不同的下列各种三核子系统态：

（7/2,1/2)，2（5/2,1/2)，（3/2,3/2)，4（3/2,1/2)，（1/2,3/2)，5（1/2,1/2).（括号前的数字代表该态的数目，$M_J$和$T_\zeta$为负值的态无需写出）.它们对应于以下各种$(J,T)$值：
\begin{equation*}
  (7/2,1/2),\,(5/2,1/2),\,(3/2,3/2),\,(3/2,1/2),\,(1/2,1/2).
\end{equation*}

\noindent{\heiti 习题3}\quad 某组态由处于$\mathrm{p}_{3/2}$态的两个中子和一个质子所组成（$n$相同），求其基态磁矩（计及同位旋不变性）\footnote{$\mathrm{Li}^7$核具有这种组态（在封闭壳层$(1\mathrm{s}_{1/2})^4$外面）.}.

解：此组态的基态具有$J=3/2$，根据本节所给规则，它的同位旋具有最小值$T=|T_\zeta|=\dfrac{1}{2}$.

现在来求对应于最大值$M_J=3/2$的系统波函数.这个值可分别来自pn，n的下列各组$m_j$值（两个核子相同时要用泡利原理）：
\begin{equation*}
  \left(\frac{3}{2},\frac{3}{2},-\frac{3}{2}\right),
  \left(\frac{3}{2},\frac{1}{2},-\frac{1}{2}\right),
  \left(\frac{1}{2},\frac{3}{2},-\frac{1}{2}\right),
  \left(-\frac{1}{2},\frac{3}{2},\frac{1}{2}\right).
\end{equation*}
因此所求的波函数$\varPsi^{JM}_{TT_\zeta}$具有下列线性组合形式：
\begin{equation}
\begin{aligned}
  \varPsi^{3/2\,3/2}_{1/2\,-1/2}
  &=a\big[\psi^{3/2}_{1/2}\psi^{3/2}_{-1/2}\psi^{-3/2}_{-1/2}\big]
  +b\big[\psi^{3/2}_{1/2}\psi^{1/2}_{-1/2}\psi^{-1/2}_{-1/2}\big]\,+\\
  &\quad+c\big[\psi^{3/2}_{-1/2}\psi^{1/2}_{1/2}\psi^{-1/2}_{-1/2}\big]
  +d\big[\psi^{3/2}_{-1/2}\psi^{1/2}_{-1/2}\psi^{-1/2}_{1/2}\big],
\end{aligned}
\end{equation}
式中的$[\cdots]$代表三个单核子波函数$\psi^{m_j}_{\tau_\zeta}$的反对称乘积（即(61.5)的行列式形式）.

下列算符作用在(1)式上必等于零（见§67，例题）：
\begin{equation*}
  \hat{T}_-=\sum_{i=1}^{3}\hat{\tau}_-^{(i)},\qquad
  J_+=\sum_{i=1}^{3}\hat{j}_+^{(i)}.
\end{equation*}
$\hat{\tau}_-^{(i)}$算符把第$i$个核子的质子波函数变成中子波函数（把中子波函数变成零）.因此很易看出，$\hat{T}_-$算符使(1)式第一项的行列式中有两行都变成零，同时把其余三项的行列式变成相等.因此得下列条件：
\begin{equation*}
  b+c+d=0.
\end{equation*}
其次，对$j=3/2$而$m_j$值不同的各个单核子态，我们有［根据(27.12)］：
\begin{equation*}
  \hat{j}_+\psi^{3/2}=0,\quad
  \hat{j}_+\psi^{1/2}=\sqrt{3}\psi^{3/2},\quad
  \hat{j}_+\psi^{-1/2}=2\psi^{1/2},\quad
  \hat{j}_+\psi^{-3/2}=\sqrt{3}\psi^{-1/2}.
\end{equation*}
由此很易找出$\hat{J}_+$算符作用在(1)式上的结果为
\begin{equation*}
  \hat{\boldsymbol{J}}_+\varPsi^{3/2\,3/2}_{1/2\,-1/2}
  =\sqrt{3}\left(a+b-c\right)
  \big[\psi^{3/2}_{1/2}\psi^{3/2}_{-1/2}\psi^{-1/2}_{-1/2}\big]
  +2\left(c-d\right)
  \big[\psi^{3/2}_{-1/2}\psi^{1/2}_{1/2}\psi^{1/2}_{-1/2}\big]
\end{equation*}
（有几项的变号与行列式中各行的置换有关）.上式等于零的条件为：
\begin{equation*}
  a+b-c=0,\qquad c-d=0.
\end{equation*}
根据以上诸条件再加上(1)式的归一化条件，可得出
\begin{equation*}
  a=\frac{3}{\sqrt{15}},\qquad
  b=-\frac{2}{\sqrt{15}},\qquad
  c=d=\frac{1}{\sqrt{15}}.
\end{equation*}
考虑到$m_j$态中质子（或中子）磁矩的平均投影值为$\mu_\mathrm{p}m_j/j$（或$\mu_\mathrm{n}m_j/j$），即可求得由(1)式波函数算出的系统磁矩平均值，它等于
\begin{equation*}
  \mu=\bar{\mu}_z=\frac{9}{15}\mu_\mathrm{p}+\frac{4}{15}\mu_\mathrm{p}
  +\frac{1}{15}\left(\frac{1}{3}\mu_\mathrm{p}+\frac{2}{3}\mu_\mathrm{n}\right)
  +\frac{1}{15}\left(-\frac{1}{3}\mu_\mathrm{p}+\frac{4}{3}\mu_\mathrm{n}\right)
  =\frac{1}{15}\left(13\mu_\mathrm{p}+2\mu_\mathrm{n}\right).
\end{equation*}
根据(118.12)，(118.13)式，处于$\mathrm{p}_{3/2}$态的核子有$\mu_\mathrm{n}=-1.91$，$\mu_\mathrm{p}=3.79$.结果得$\mu=3.03$.

\noindent{\heiti 习题4}\quad 设满壳层外的所有核子都属于同一态，并且质子数等于中子数，求核磁矩.

解：由于$N=Z$时同位旋分量值$T_\zeta=0$，因此只有下列算符的同位标量部分才有对角矩阵元：
\begin{equation*}
  \hat{\boldsymbol{\mu}}=\sum_n g_\mathrm{n}\hat{\boldsymbol{j}}_n+\sum_p g_\mathrm{p}\hat{\boldsymbol{j}}_p;
\end{equation*}
参考§116之末.根据(116.5)式取出以上算符的同位标量部分，它等于
\begin{equation*}
  \frac{1}{2}\left(g_\mathrm{n}+g_\mathrm{p}\right)\sum_{n,p}\hat{\boldsymbol{j}}
  =\frac{1}{2}\left(g_\mathrm{n}+g_\mathrm{p}\right)\hat{\boldsymbol{J}}.
\end{equation*}
因此总平均核磁矩等于$\dfrac{1}{2}\left(g_\mathrm{n}+g_\mathrm{p}\right)J$.

\noindent{\heiti 习题5}\quad 计算角动量为$j$的核子的附加磁矩，用(118.6)式的自旋-轨道分裂值把它表达出来（M.~Goepper-Mayer, J.~Jensen,1952）.

解：对算符(118.14)的角部求平均［把(118.14)大括号内的表式记作$\hat{\boldsymbol{\sigma}}$］.应用§29例题中所得的公式，得出结果为
\begin{equation}
  \bar{\boldsymbol{\sigma}}\equiv\overline{\hat{\boldsymbol{s}}-(\hat{\boldsymbol{s}}\cdot\boldsymbol{n})\boldsymbol{n}}
  =\frac{2}{3}\hat{\boldsymbol{s}}
  +\frac{(\hat{\boldsymbol{s}}\cdot\hat{\boldsymbol{l}})\hat{\boldsymbol{l}}
  +(\hat{\boldsymbol{l}}\cdot\hat{\boldsymbol{s}})\hat{\boldsymbol{l}}
  -\dfrac{2}{3}l(l+1)\hat{\boldsymbol{s}}}{(2l-1)(2l+3)}
\end{equation}
另一方面，对核子运动整个平均以后，$\boldsymbol{\sigma}$的平均值只能沿$j$方向，即$\bar{\boldsymbol{\sigma}}=a\hat{\boldsymbol{j}}$；其中的$a=(\bar{\boldsymbol{\sigma}}\cdot\hat{\boldsymbol{j}})/j^2$.把(2)式中的矢量投影到$j$上［并且考虑到算符$\hat{\boldsymbol{j}}$和$(\hat{\boldsymbol{l}}\cdot\hat{\boldsymbol{s}})$是对易的］同时把$\boldsymbol{l}\cdot\boldsymbol{s}$，$\boldsymbol{l}^2$等等改成它们的本征值，经过简单计算后可得下式所示的核子附加磁矩（以核磁子为单位）：
\begin{equation}
  \mu_{\text{附加}}=\mp\overline{f(r)}\,\frac{m_\mathrm{p}R^2}{\hbar^2}\,
  \frac{2j+1}{4(j+1)}\quad\left(j=l\pm\frac{1}{2}\text{时}\right),
\end{equation}
$m_\mathrm{p}$是核子质量，$R$是核半径.由于$f(r)$在核内深处很快地衰减，求$r^2f$的平均时可以把$r^2$改成$R^2$.(3)中的$\bar{f}$可按(118.6)式表为自旋-轨道分裂值.

\end{xiti}
\renewcommand{\theequation}{\arabic{section}.\arabic{equation}}%

\section{非球形核}

在有心力场中运动的粒子组不可能具有转动能谱，在量子力学中，这个系统的转动概念是没有意义的.这一点适用于§117中所讲的具有球对称自洽场的原子核壳层模型.

把一个系统的能量分成“内在”和“转动”部分，这在量子力学中没有确切的含义，它只能是近似的并且在下列情形下才是可能的，即由于某种物理原因，这个系统能够很好地近似成为运动于某一给定非球对称场中的粒子组.考虑到这种场相对于某一固定坐标系转动的可能性，结果就有能级的转动结构.例如分子中就出现这样的情形，它的电子谱项可以作为运动于给定的固定核场中的一个多电子系统的能级来加以确定.

实验表明，大多数原子核确实没有转动结构，这意味着对这些核来说，球对称自洽场是一个很好的近似.也就是说，除有量子涨落外，这些核是球形的.

但也存在另一类核，它们具有转动型的能谱；它们大致位于原子量为$150<A<190$和$A>220$的范围内，这个性质意味着，球对称的自洽场差不多对这些核是完全不能适用的，而从原则上讲，应在事先不作对称性假定的情形下去求自洽场，以便使这些核的形状也能“自洽地”加以确定.实验表明，这类原子核的正确模型是这样的，它的自洽场具有一个对称轴和一个垂直于它的对称平面（亦即具有旋转椭球的对称性）.非球形核的概念在A.~Bohr和B.~R.~Mottelson（1952---1953）的工作中得到了详尽的发展.

需要强调的是，在这里我们考虑的是性质上不同的两类核.这一点特别可从下列事实看出，核或者是球形的，或者就是“形变程度”并不很小的非球形.

核内未满壳层的存在有利于非球形的出现，而核子的配对现象看来也相当重要.另一方面，封闭壳层趋向于给出球形核，一个典型例子是双幻核$^{208}_{82}\mathrm{Pb}$：由于它的核子组态的封闭性，这个核（以及邻近于它的核）是球形的，这就使得在非球形重核序列中出现了一个间断.

非球形核的能级由两部分组成：“固定”核的能级和整体转动的能级.在偶-偶核中，转动结构的能级间距小于“固定”核的能级间距.

非球形核的能级分类在许多方面类似于由两个相同原子组成的双原子分子，因为这两种情形中粒子（核子或电子）所处的场具有相同的对称性.因此我们可以直接应用第十一章中所得的一系列结论\footnote{应该指出，我们所讲的能级分类的相似性指的是双原子分子不是指对称陀螺.一个多粒子系统在轴对称的场中运动，绕场轴转动的概念就失去意义，正像有心力场中的系统绕任一轴旋转的概念一样.}.

我们先考虑“固定”核的状态分类，在轴对称场中，只有角动量沿对称轴的分量是守恒量.因此每个核态首先由总角动量的$\Omega$分量所描写\footnote{按定义有$\Omega\geqslant0$（正如双原子分子中的量子数$\Lambda$是正的一样），记得在双原子分子中，仅当$\Omega$被定义成为$\Lambda+\Sigma$并且$\Sigma$可正可负（取决于轨道角动量和自旋的相对方向）时，$\Omega$才能具有负值.}，它可以是整数或半整数.根据波函数在所有核子坐标（相对于该核中心）反号时的行为，能级可用偶$(g)$或奇$(u)$描述.

此外，$\Omega=0$时，根据波函数对通过核轴的平面的反射行为，可分为正态和负态（见§78）.

偶-偶非球形核的基态为$0_g^+$（零代表$\Omega$值），对应于具有零角动量和最高对称性的波函数.这是所有中子和质子配对的结果.如果原子核含有奇数个质子或中子，我们可以考虑在偶-偶核实的自洽场中的“单个”核子态.此时的$\Omega$值由该核子的角动量分量$\omega$所确定.同理，奇-奇核中的$\Omega$值由奇数中子和质子的角动量分量所确定：
\begin{equation*}
  \Omega=|\omega_\mathrm{p}\pm\omega_\mathrm{n}|.
\end{equation*}
需要强调的是，我们不能说核子的自旋分量及其轨道角动量分量同时具有定值.理由是这样的，虽然核子的自旋轨道耦合小于它和核实自洽场的作用能，但它并不小于假如微扰论可用（从而核子的自旋和角动量可以近似地分开考虑\footnote{然而在球形核中却有定义量$l$的可能，这是由于宇称和角动量的同时守恒.}）时核子在自洽场中应有的相邻能级间距.

现在考虑非球形核的转动结构.这个结构中的间距小于原子核中核子的自旋-轨道作用.这相当于双原子分子理论中的情形$a$（§83）.

转动核的总角动量$\boldsymbol{J}$当然是守恒的，对给定的$\Omega$，$J$的取值是从$\Omega$开始：
\begin{equation}
  J=\Omega,\ \Omega+1,\ \Omega+2,\cdots;
\end{equation}
见(83.2).对$\Omega=0$的原子核，$J$的可能值还有一个附加限制：$0_g^+$和$0_u^-$态中的$J$值只能取偶数，而$0_g^-$和$0_u^+$态中的只能取奇数（见§86）.特别是偶-偶核基项（$0_g^+$）的转动能级中，$J$的取值为$0,2,4,\cdots$.

原子核的转动能量由下式给出：
\begin{equation}
  E_{\text{转动}}=\frac{\hbar^2}{2I}J(J+1),
\end{equation}
$I$是原子核的转动惯量（绕垂直于对称轴的一个轴）；这个公式对应于双原子分子理论中的类似表式［(83.6)式中依赖于$J$的那一项］.最低能级对应于$J$的最低值，即$J=\Omega$.

根据(119.2)，能级的转动结构可用某些间隔规则来描写，这些规则和能级（$\Omega$为给定）的其它特征无关.以偶-偶核基项的转动结构部分（具有$J=2,4,6,8,\cdots$）为例，从最低能级（$J=0$）开始各间距的比例为$1:3.3:7:12\cdots$.

但是，(119.2)对$\Omega=\dfrac{1}{2}$的态讲来是不够的，这种态可以在原子核中有奇数个核子时出现.这种情形下，还有一项与(119.2)差不多大小的能量贡献，它来自核子和转动核离心场的相互作用.此项与$J$的关系可用下法找出.

力学中知道（《力学》§39），转动坐标系中的粒子能量含有一个附加项，它等于转动角速度和粒子角动量的乘积.原子核哈密顿算符中的这一相应项可以写成$2b\hat{\boldsymbol{K}}\cdot\hat{\boldsymbol{\sigma}}$的形式，式中$b$是某个常数，$\boldsymbol{K}$是核实（去掉一个外核子后的原子核）的角动量，$\boldsymbol{\sigma}$是该核子的角动量.后者必须从纯形式意义来理解，实际上，在原子核的轴对称场中并不存在核子的角动量矢量，把$\hat{\boldsymbol{\sigma}}$看作类似于自旋为$\dfrac{1}{2}$算符的含义是，它能给出角动量分量值为$\pm\dfrac{1}{2}$的两态之间的跃迁，与$\Omega=\dfrac{1}{2}$的情形相一致\footnote{$\Omega=\frac{1}{2}$情形的特点是，只有属于同一能级并且角动量分量值反号的两个态之间才存在着能量微扰项的跃迁矩阵元，这就使得即使是微扰论的一级近似中也出现能级位移.
这种现象与$\Omega=\frac{1}{2}$的双原子分子能级的$\Lambda$双线（§88）相类似.}.由于$\boldsymbol{K}=\boldsymbol{J}-\boldsymbol{\sigma}$，这个算符的本征值为
\begin{equation*}
  2b\boldsymbol{K}\cdot\boldsymbol{\sigma}
  =b\left[J(J+1)-K(K+1)-\frac{3}{4}\right].
\end{equation*}
为方便计，我们在上式中加进一项与$J$无关的常数项$b$，当$J=K\pm\dfrac{1}{2}$时上式就等于$\pm b\left(J+\dfrac{1}{2}\right)$.

如果利用偶-偶核实的角动量$K$是一个偶数的事实，上述表式可写成$(-1)^{J-\frac{1}{2}}b\left(J+\dfrac{1}{2}\right)$.因此对$\Omega=\dfrac{1}{2}$的原子核转动能量，可得下列表示式：
\begin{equation}
  E_{\text{转动}}=\frac{\hbar^2}{2I}J(J+1)
  +(-1)^{J-\frac{1}{2}}b\left(J+\frac{1}{2}\right).
\end{equation}
（A.~Bohr, B.~R.~Mottelson,1953）.注意当常数$b$是正的并且足够大时，$J=3/2$的能级可能处于$J=\dfrac{1}{2}$能级的下面，亦即转动能级的正常次序（最低能级对应于$J$的最小允许值）将会改变.

非球形核的转动惯量不能像具有给定形状的刚体那样去计算.这样的计算只有当运动于核自洽场中的核子可以看作彼此间没有直接作用时才有可能.实际上，成对效应导致转动惯量的减小，使它小于刚体值.

非球形核的磁矩$\boldsymbol{\mu}$包括“固定”核的磁矩和来自核转动的磁矩，前者（对核内的核子运动平均以后）沿核轴，可写成$\mu'\boldsymbol{n}$，$\mu'$是它的值，$\boldsymbol{n}$是沿核轴的单位矢量，来自转动的磁矩（经过同样的平均以后）沿矢量$\boldsymbol{J}-\Omega\boldsymbol{n}$方向，这个矢量是原子核的总角动量减去“固定核”中核子的总角动量\footnote{这个写法仅适用于$\Omega\neq\frac{1}{2}$的情形（见题2）.}.故
\begin{equation}
  \boldsymbol{\mu}=\mu'\boldsymbol{n}+g_r\left(\boldsymbol{J}-\Omega\boldsymbol{n}\right).
\end{equation}
$g_r$是转动核的旋磁因子.由于转动中的磁矩仅由质子所贡献，我们有
\begin{equation}
  g_r=\frac{I_\mathrm{p}}{I_\mathrm{p}+I_\mathrm{n}},
\end{equation}
式中$I_\mathrm{n}$和$I_\mathrm{p}$是原子核转动惯量的中子部分和质子部分.对一个质子系统，简单地有$g_r=1$.一般说来，(119.5)式的比值不等于原子核的质子数和质量数之比$Z/A$.

对核转动平均以后，磁矩沿守恒矢量$\boldsymbol{J}$的方向：
\begin{equation*}
  \hat{\boldsymbol{\mu}}=\frac{\mu}{J}\hat{\boldsymbol{J}}
  =\left(\mu'-\Omega g_r\right)\bar{\boldsymbol{n}}+g_r\hat{\boldsymbol{J}}.
\end{equation*}
和通常做法一样，上式两边乘$\hat{\boldsymbol{J}}$后取本征值，对$\Omega=J$的基态核，有
\begin{equation}
  \mu=\left(\mu'+g_r\right)\frac{J}{J+1}.
\end{equation}

\begin{xiti}

\noindent{\heiti 习题1}\quad 试用$Q_0$表达转动核的四极矩$Q$，$Q_0$为相对于核轴（固定在核上）的四极矩（A.~Bohr 1951\mbox{）.}

解：转动核四极矩张量算符可通过$Q_0$表成
\begin{equation*}
  Q_{ik}=\frac{3}{2}Q_0\left(n_in_k-\frac{1}{3}\delta_{ik}\right).
\end{equation*}
这是一个零迹对称张量，由核轴单位矢量$\boldsymbol{n}$的分量所组成，并且$Q_{zz}=Q_0$.对原子核转动态求平均的方法类似于§29中例题之解（不同之处是$n_i\overline{J}_i=\Omega$，而不是零），从而得到(75.2)那样的表式，而
\begin{equation*}
  Q=Q_0\,\frac{3\Omega^2-J(J+1)}{(2J+3)(J+1)}.
\end{equation*}
对$\Omega=J$的原子核基态，我们得
\begin{equation*}
  Q=Q_0\,\frac{(2J-1)J}{(2J+3)(J+1)}
\end{equation*}
$J$增大时比值$Q/Q_0$趋于1，但是很慢.

\noindent{\heiti 习题2}\quad 求$\Omega=\dfrac{1}{2}$的原子核基态磁矩.

解：此时的磁矩算符可用本节引进的$\hat{\boldsymbol{\sigma}}$算符写成下列形式
\begin{equation*}
  \hat{\boldsymbol{\mu}}=2\mu'\hat{\boldsymbol{\sigma}}+g_r\hat{\boldsymbol{K}},\qquad
  \hat{\boldsymbol{K}}=\hat{\boldsymbol{J}}-\hat{\boldsymbol{\sigma}}.
\end{equation*}
以下的计算和本节介绍的相同.如果$J=\dfrac{1}{2}$对应于原子核的基态（$K=J-\dfrac{1}{2}=0$），我们有$\mu=\mu'$；如果基态中$J=3/2$（以及$K=J+\dfrac{1}{2}=2$），则$\mu=(9g_r-3\mu')/5$.

\noindent{\heiti 习题3}\quad 求偶-偶核基态转动结构的前几个能级的能量，该核具有旋转椭球对称性.

解：偶-偶核基态对应于最对称的“固定”核波函数，亦即具有$\boldsymbol{D}_2$群的$A$表示对称性的波函数.因此对给定的$J$值共有$\dfrac{1}{2}J+1$个（$J$为偶数时）或$\dfrac{1}{2}(J-1)$个（$J$为奇数时）不同的能级.$J=2$时由§103题3的(7)式给出，$J=3$时由§103题4的(8)式给出.

\end{xiti}

\section{同位素移位}

原子核的特殊性质（有限的质量，大小，自旋）使它不同于库仑场有固定力心，这对原子的电子能级产生一定的影响.

其中的一个效应称为能级的同位素移位，这是原子从一个同位素变到另一个同位素时的能级变动.当然，我们感兴趣的实际上并不是一个能级的变动，而是谱线中观测到的能级间距的变动，由于这一原因，我们实际上需要考虑的并不是原子的整个电子壳层的能量，而只是与参与跃迁的电子有关的那一部分能量.

轻原子中，同位素移位主要来自核的有限质量，计及核的运动后，哈密顿算符中出现以下一项：
\begin{equation*}
  \frac{1}{2M}\left(\sum_i\hat{\boldsymbol{p}}_i\right)^2,
\end{equation*}
式中的$M$是核的质量，$\boldsymbol{p}_i$是电子的动量\footnote{在原子的质心系中，原子核动量和电子动量之和等于零：$\boldsymbol{p}_{\text{核}}+\sum\boldsymbol{p}_i=0$. 因此它们的总动能为
\begin{equation*}
  \frac{\boldsymbol{p}_{\text{核}}^2}{2M}
  +\frac{1}{2m}\sum\boldsymbol{p}_i^2
  =\frac{1}{2M}\biggl(\sum_i\boldsymbol{p}_i\biggr)^2
  +\frac{1}{2m}\sum\boldsymbol{p}_i^2.
\end{equation*}}.来自这个效应的同位素移位因而由以下平均值给出：
\begin{equation}
  \frac{1}{2}\left(\frac{1}{M_1}-\frac{1}{M_2}\right)
  \overline{\left(\sum_i\boldsymbol{p}_i\right)^2},
\end{equation}
这个平均是由有关原子态的波函数算出的（$M_1$和$M_2$是两个同位素的核质量）.

重原子中，同位素移位的主要贡献来自原子核的有限大小.这个效应实际上只对处于s态的外电子能级才是显著的，由于s态波函数（与$l\neq0$的态不同）当$r\to0$时不趋于零，所以在“核内”找到电子的概率比较大.我们来计算这种情形下的同位素移位\footnote{以下给出的计算中没有考虑到原子核附近电子运动的相对论效应，从而只当条件$Ze^2/\hbar c\ll1$满足时才能成立.}.

令$\varphi(r)$为核场的实际静电势，不同于点电荷$Ze$的库仑势$Ze/r$.与纯库仑场$Ze/r$中的值相比较，电子能量的改变由下列积分给出：
\begin{equation}
  \Delta E=-e\int\left(\varphi-\frac{Ze}{r}\right)\psi^2(r)\,\mathrm{d}V,
\end{equation}
式中的$\psi(r)$是电子波函数，s态中这个函数是球对称的并且是实函数.上式中的积分从形式上讲来虽然延及整个空间，但实际上被积函数中的$\varphi-Ze/r$只在核体积内才不等于零.另一方面，当$r\to0$时s态波函数趋于一个常数（参考§32），而且实际上甚至在核外就已经达到了这个常数值.因此可把$\psi^2$移出积分号外，并把这个由库仑点电荷算出的$\psi(r)$改成它在$r=0$处的值.

为了进一步变换这个积分，利用恒等式$\Delta r^2=6$并把(120.2)写成以下形式：
\begin{equation*}
  \Delta E=-\frac{1}{6}e\psi^2(0)\int\left(\varphi-\frac{Ze}{r}\right)\Delta r^2\,\mathrm{d}V
  =-\frac{1}{6}e\psi^2(0)\int r^2\Delta\left(\varphi-\frac{Ze}{r}\right)\mathrm{d}V.
\end{equation*}
体积分的变换中已经利用了无穷远面上的积分等于零的事实.但是$\Delta\dfrac{1}{r}=-4\pi\delta(\boldsymbol{r})$，而且对所有的$\boldsymbol{r}$有$r^2\delta(\boldsymbol{r})=0$.再根据静电学的泊松公式$\Delta\varphi=-4\pi\rho$，式中的$\rho$现在是原子核内的电荷密度分布.最后得下列结果：
\begin{equation}
  \Delta E=\frac{2\pi}{3}\psi^2(0)Ze^2\,\overline{r^2},
\end{equation}
式中
\begin{equation*}
  \overline{r^2}=\frac{1}{Ze}\int\rho r^2\,\mathrm{d}V
\end{equation*}
是原子核的质子均方半径.核内质子均匀分布时，$\overline{r^2}=3R^2/5$，$R$是核的几何半径.能级的同位素移位等于两个同位素的(120.3)式之差.

§71中曾经估计过$\psi(0)$，表明它和原子序（假定很大）的关系为$\sqrt{Z}$.因此(120.3)所表示的分裂值和$R^2Z^2$成正比.

\section{原子能级的超精细结构}

原子中来自原子核性质的另一个效应，是电子和核自旋相互作用引起的原子能级分裂，这称为能级的超精细结构.考虑到这种作用很弱，超精细结构的间距即使和精细结构的间距相比也小得很多，所以超精细结构必须针对每个精细结构成分分别加以考虑.

本节中用$i$代表核自旋（与原子光谱中的常用记号一致），记号$J$仍留作原子中电子壳层的总角动量.原子（包括原子核）的总角动量记作$\boldsymbol{F}=\boldsymbol{J}+\boldsymbol{i}$.每个超精细结构成分由这个角动量的某一个定值所描写.按照角动量相加的一般规则，量子数$F$的取值为
\begin{equation}
  F=J+i,\ J+i-1,\cdots,\ |J-i|,
\end{equation}
因此每个具有给定$J$值的能级分裂成为$2i+1$个（如果$i<J$）或者$2J+1$个（如果$i>J$）能级.

由于原子中电子间的平均距离$r$远大于核半径$R$，电子和最低阶核多极矩的作用在超精细分裂中起着重要的作用.这些矩计有磁偶极矩和电四极矩，平均电偶极矩则为零（见§75）.

核磁矩的数量级为$\mu_{\text{核}}\sim eRv_{\text{核}}/c$，其中$v_{\text{核}}$是原子核中核子的速度.它和电子磁矩（$\mu_{\text{电子}}\sim e\hbar/mc$）的相互作用能量具有数量级
\begin{equation}
  \frac{\mu_{\text{核}}\mu_{\text{电子}}}{r}
  \sim\frac{e^2\hbar}{mc^2}\,\frac{Rv_{\text{核}}}{r^3}.
\end{equation}
四极矩$Q\sim eR^2$，它所产生的场与电子电荷的相互作用能量具有数量级
\begin{equation}
  \frac{eQ}{r^3}\sim\frac{e^2R^2}{r^3}.
\end{equation}
(121.2)和(121.3)的比较表明，磁作用（及其产生的能级分裂）要比四极矩作用大$\left(v_{\text{核}}/c\right)\left(\hbar/mcR\right)\sim15$倍，尽管比值$v_{\text{核}}/c$比较小，但比值$\hbar/mcR$是大的.

电子和原子核的磁作用算符具有下列形式：
\begin{equation}
  \hat{V}_{iJ}=a\hat{\boldsymbol{i}}\cdot\hat{\boldsymbol{J}}
\end{equation}
［类似于电子的自旋-轨道作用(72.4)］.它所产生的能级分裂和$F$的关系因此为
\begin{equation}
  \frac{1}{2}aF(F+1);
\end{equation}
参考(72.5).

电子和核四极矩的相互作用算符，是由核四极矩张量算符$\hat{Q}_{ik}$和电子的角动量矢量$\hat{\boldsymbol{J}}$的分量构成的.它与这些算符组成的标量$\hat{Q}_{ik}\hat{J}_i\hat{J}_k$成正比，具有下列形式：
\begin{equation}
  b\left[\hat{i}_i\hat{i}_k+\hat{i}_k\hat{i}_i
  -\frac{2}{3}i(i+1)\delta_{ik}\right]\hat{J}_i\hat{J}_k;
\end{equation}
式中应用了这样的事实，即$Q_{ik}$通过(75.2)那样的公式用核自旋算符表出.算出算符(121.6)的本征值（与§84中题1的算法完全类似），可得能级的四极矩超精细分裂与量子数$F$的下列关系式：
\begin{equation}
  \frac{1}{2}bF^2(F+1)^2
  +\frac{1}{2}bF(F+1)\left[1-2J(J+1)-2i(i+1)\right].
\end{equation}
磁偶极超精细分裂效应对外层的s态电子的能级特别显著，因为这样的一个电子有较大的概率处于原子核附近.

我们来计算含有一个外层s电子的原子的超精细分裂（E.~Fermi,1930）.这个电子由球对称的$\psi(r)$波函数所描写，它在其余电子和原子核的自洽场中运动\footnote{以下的计算假定满足条件$Ze^2/\hbar c\ll1$（见 452 页注{\symbol{"2460}}）.}.

我们把电子和核的相互作用算符取作$-\hat{\boldsymbol{\mu}}\cdot\hat{\boldsymbol{H}}$，这是核磁矩$\hat{\boldsymbol{\mu}}=\mu\hat{\boldsymbol{i}}/i$在电子所产生的磁场$\boldsymbol{H}$（原点处的磁场）中的能量算符.根据电动力学中的熟知公式，此磁场为
\begin{equation}
  \hat{\boldsymbol{H}}=\frac{1}{c}\int\frac{\boldsymbol{n}\times\hat{\boldsymbol{j}}}{r^2}\,\mathrm{d}V.
\end{equation}
式中的$\hat{\boldsymbol{j}}$是来自运动电子的自旋的电流密度算符，$\boldsymbol{r}=r\boldsymbol{n}$是原点到体积元$\mathrm{d}V$的径矢\footnote{见《场论》(43.7). 该式中的矢量$\boldsymbol{R}$是反方向的，即从$\mathrm{d}V$到原点（场的观测点）.}.按照(115.4)，
\begin{equation*}
  \hat{\boldsymbol{j}}=-2\mu_\mathrm{B}c\,\mathrm{curl}(\psi^2\hat{\boldsymbol{s}})
  =-2\mu_\mathrm{B}c\,\frac{\mathrm{d}\psi^2(r)}{\mathrm{d}r}\boldsymbol{n}\times\hat{\boldsymbol{s}},
\end{equation*}
$\mu_\mathrm{B}$是玻尔磁子.令$\mathrm{d}V=r^2\mathrm{d}r\,\mathrm{d}o$并进行积分，得
\begin{equation*}
  \hat{\boldsymbol{H}}=-2\mu_\mathrm{B}\int_0^{\infty}\frac{\mathrm{d}\psi^2}{\mathrm{d}r}\mathrm{d}r
  \int\boldsymbol{n}\times(\boldsymbol{n}\times\boldsymbol{s})\,\mathrm{d}o
  =-2\mu_\mathrm{B}\psi^2(0)\frac{8\pi}{3}\hat{\boldsymbol{s}}.
\end{equation*}
相互作用算符最后变成
\begin{equation}
  \hat{V}_{is}=-\hat{\boldsymbol{\mu}}\cdot\hat{\boldsymbol{H}}
  =\frac{16\pi}{3i}\mu\mu_\mathrm{B}\psi^2(0)\hat{\boldsymbol{i}}\cdot\hat{\boldsymbol{s}}.
\end{equation}
如果原子的总角动量$J=S=\dfrac{1}{2}$，超精细分裂产生的是双线（$F=i\pm\dfrac{1}{2}$）；根据(121.5)和(121.9)，这两个能级的间距为
\begin{equation}
  E_{i+\frac{1}{2}}-E_{i-\frac{1}{2}}
  =\frac{8\pi}{3i}\mu\mu_\mathrm{B}\left(2i+1\right)\psi^2(0).
\end{equation}
由于$\psi(0)$值正比于$\sqrt{Z}$（见§71），这个分裂值正比于原子序.

\begin{xiti}

\noindent{\heiti 习题1}\quad 求原子的超精细分裂（来自磁作用），该原子具有轨道角动量为$l$的一个（封闭壳层外的）电子（E.~Fermi,1930）.

解：核磁矩$\boldsymbol{\mu}$产生的矢势和场强为
\begin{equation*}
  \boldsymbol{A}=\frac{\boldsymbol{\mu}\times\boldsymbol{n}}{r^2},\qquad
  \boldsymbol{H}=\frac{3\boldsymbol{n}(\boldsymbol{\mu}\cdot\boldsymbol{n})-\boldsymbol{\mu}}{r^3}
\end{equation*}
（$\mathrm{div}\,\boldsymbol{A}=0$）.应用这些式子，可把相互作用算符写成下列形式：
\begin{equation*}
  \frac{|e|}{mc}\boldsymbol{A}\cdot\hat{\boldsymbol{p}}
  +\frac{|e|\hbar}{mc}\hat{\boldsymbol{H}}\cdot\hat{\boldsymbol{s}}
  =\frac{2\mu_\mathrm{B}}{r^3}\hat{\boldsymbol{\mu}}\cdot
  \left[\hat{\boldsymbol{l}}+3(\hat{\boldsymbol{s}}\cdot\boldsymbol{n})\boldsymbol{n}-\hat{\boldsymbol{s}}\right].
\end{equation*}
对具有给定$j$值的态平均以后，方括号内的表式沿$\boldsymbol{j}$方向.因而可写出
\begin{equation*}
  \hat{V}_{ij}=2\mu_\mathrm{B}\hat{\boldsymbol{\mu}}\cdot\hat{\boldsymbol{j}}
  \left[\hat{\boldsymbol{l}}\cdot\hat{\boldsymbol{j}}
  +3(\hat{\boldsymbol{s}}\cdot\bar{\boldsymbol{n}})(\bar{\boldsymbol{n}}\cdot\hat{\boldsymbol{j}})
  -\hat{\boldsymbol{s}}\cdot\hat{\boldsymbol{j}}\right]
  \frac{\overline{r^{-3}}}{j(j+1)},
\end{equation*}
$n_in_k$的平均值已在§29例题中算出.采用它并取本征值后得
\begin{equation*}
  \frac{2\mu_\mathrm{B}\mu}{i}\,\hat{\boldsymbol{i}}\cdot\hat{\boldsymbol{j}}
  \left[\boldsymbol{l}\cdot\boldsymbol{j}
  +\frac{2l(l+1)\boldsymbol{s}\cdot\boldsymbol{j}-6(\boldsymbol{s}\cdot\boldsymbol{l})(\boldsymbol{j}\cdot\boldsymbol{l})}{(2l-1)(2l+3)}\right]
  \frac{\overline{r^{-3}}}{j(j+1)},
\end{equation*}
再经简单计算，最后得
\begin{equation*}
  \frac{\mu_\mathrm{B}\mu}{i}\,\frac{l(l+1)}{j(j+1)}F(F+1)\,\overline{r^{-3}},
\end{equation*}
其中$\boldsymbol{F}=\boldsymbol{j}+\boldsymbol{i}$和$j=l\pm\dfrac{1}{2}$.$\overline{r^{-3}}$的平均是对电子波函数的径向部分求的.

\noindent{\heiti 习题2}\quad 求原子能级超精细结构分量的塞曼分裂（S.~A.~Goudsmit, R.~F.~Bacher, 1930）.

解：(113.4)式中（假定外场很弱，它所产生的分裂小于超精细结构间距），现在不但要对电子态求平均并且还要对核自旋的方向求平均.从第一个平均得$\Delta E=\mu_\mathrm{B}g_JJ_zH$，其中$g_J$同以前的(113.7)式一样.第二个平均给出［类似于(113.5)］：
\begin{equation*}
  \bar{J}_z=(\boldsymbol{J}\cdot\boldsymbol{F})M_F/F^2,
\end{equation*}
故最后得
\begin{equation*}
  \Delta E=\mu_\mathrm{B}g_FHM_F,\qquad
  g_F=g_J\frac{F(F+1)+J(J+1)-i(i+1)}{2F(F+1)}.
\end{equation*}

\end{xiti}

\section{分子能级的超精细结构}

分子能级的超精细结构与原子能级的相类似.

大多数分子中电子总自旋为零.能级超精细分裂的主要来源就是电子和核的四极矩作用.当然，只有自旋$i$不等于0和$\dfrac{1}{2}$的那些核才能参与这种作用，否则它们的四极矩等于零.

鉴于分子中的核运动比较慢，四极矩相互作用算符对分子态的平均可以分成两步：先对固定核的电子态求平均，再对分子的转动求平均.

我们先考虑双原子分子.第一步平均给出电子和每个核的作用，可表成正比于标量$\hat{Q}_{ik}n_in_k$的一个算符，由核四极矩张量算符和分子轴的单位矢量$\boldsymbol{n}$所组成，$\boldsymbol{n}$是确定该分子与核自旋相对取向的唯一矢量.由于$\hat{Q}_{ii}=0$，这个算符可写成下列形式：
\begin{equation}
  b\hat{i}_i\hat{i}_k\left(n_in_k-\frac{1}{3}\delta_{ik}\right).
\end{equation}
给定了核自旋沿分子轴的分量$i_\zeta$后，这个量等于$b\left[i_\zeta^2-\dfrac{1}{3}i(i+1)\right]$.

当算符(122.1)对分子转动求平均后，它可通过守恒的转动角动量算符$\hat{\boldsymbol{K}}$表达出来.乘积$n_in_k$的平均可用§29例题中导出的公式（矢量$\boldsymbol{l}$改成$\boldsymbol{K}$），其结果为
\begin{equation}
  -\frac{b}{(2K-1)(2K+3)}\hat{i}_i\hat{i}_k
  \left[\hat{K}_i\hat{K}_k+\hat{K}_k\hat{K}_i
  -\frac{2}{3}\delta_{ik}K(K+1)\right].
\end{equation}
这个算符的本征值可用(121.6)中的同样做法求出.

对于多原子分子，(122.1)一般地改成下列形式的一个算符：
\begin{equation}
  b_{ik}\hat{i}_i\hat{i}_k,
\end{equation}
$b_{ik}$是标志该分子电子态的一个零迹张量.对分子转动平均以后，这个张量即可通过总转动角动量$\boldsymbol{J}$表成下列形式：
\begin{equation}
  \bar{b}_{ik}=b\left[\hat{J}_i\hat{J}_k+\hat{J}_k\hat{J}_i
  -\frac{2}{3}J(J+1)\delta_{ik}\right].
\end{equation}
系数$b$原则上可用$b_{ik}$张量沿惯量主轴$\xi$，$\eta$，$\zeta$的三个分量表达出来；由于这些轴和分子固定在一起，$b_{\xi\xi}$等分量作为分子的一种性质不受平均的影响.我们来考虑标量$\overline{b_{ik}J_iJ_k}$，用(122.4)计算得
\begin{equation}
  \overline{b_{ik}J_iJ_k}=bJ(J+1)\left[\frac{4}{3}J(J+1)-1\right];
\end{equation}
算法类似于§29中的例题.把张量乘积展成沿$\xi$，$\eta$，$\zeta$轴的分量，我们得
\begin{equation}
  \overline{b_{ik}J_iJ_k}
  =b_{\xi\xi}\overline{J}_\xi^2+b_{\eta\eta}\overline{J}_\eta^2
  +b_{\zeta\zeta}\overline{J}_\zeta^2,
\end{equation}
这里应用了乘积$J_\xi J_\zeta$等等的平均值等于零这一事实\footnote{在$\boldsymbol{J}$的一个分量（例如$J_\zeta$）为对角的矩阵表示中，乘积$J_\xi J_\zeta$，$J_\eta J_\zeta$等等只有量子数$k$改变 1 的那些跃迁矩阵元才不等于零，可是非对称陀螺定态波函数所含的$\psi_{Jk}$中$k$的差值是一个偶数（见§103）.}.$\overline{J}_\xi^2$等等的平均值原则上可用陀螺的相应转动态波函数求出.特别是对称陀螺，简单地有
\begin{equation*}
  \overline{J}_\zeta^2=k^2,\qquad
  \overline{J}_\xi^2=\overline{J}_\eta^2
  =\frac{1}{2}\left[J(J+1)-k^2\right].
\end{equation*}
如果核自旋为$\dfrac{1}{2}$，四极矩作用不存在.在此情形下超精细分裂的一个主要来源是核磁矩之间的直接磁作用.两个磁矩$\boldsymbol{\mu}_1=\mu_1\boldsymbol{i}_1/i_1$和$\boldsymbol{\mu}_2=\mu_2\boldsymbol{i}_2/i_2$间的相互作用算符为
\begin{equation*}
  \frac{\mu_1\mu_2}{i_1i_2r^3}
  \left[\hat{\boldsymbol{i}}_1\cdot\hat{\boldsymbol{i}}_2
  -3(\boldsymbol{i}_1\cdot\boldsymbol{n})(\boldsymbol{i}_2\cdot\boldsymbol{n})\right].
\end{equation*}
计算分裂能量时，如前所述，必须对分子态求平均.

当分子中含有重原子时，核磁矩间的直接作用和通过电子壳层的间接作用对超精细分裂都有相当大的贡献.从形式上讲，这种相互作用相对于核自旋与电子的作用讲来只属于微扰论的二级近似效应.根据§121的结论很易看出，这个效应和核矩直接作用的比值为$\left(Ze^2/\hbar c\right)^2$的数量级，当$Z$很大时，它接近于1.

最后，分子能级超精细分裂中某些贡献来自核磁矩与分子转动的作用.转动分子是一个电荷的运动系统，产生一定的磁场，给出了电流密度$\boldsymbol{j}=\rho\boldsymbol{\Omega}\times\boldsymbol{r}$后即可用电动力学的公式算出磁场，上式中的$\rho$是静止分子的（电子和核的）电荷密度，$\boldsymbol{\Omega}$是转动角速度.能级分裂值可从这个磁场中的核磁矩能量算出，分子的角速度分量应该通过它的角动量分量表达出来（参考§103）.
